% Options for packages loaded elsewhere
\PassOptionsToPackage{unicode}{hyperref}
\PassOptionsToPackage{hyphens}{url}
\documentclass[
]{article}
\usepackage{xcolor}
\usepackage[margin=1in]{geometry}
\usepackage{amsmath,amssymb}
\setcounter{secnumdepth}{-\maxdimen} % remove section numbering
\usepackage{iftex}
\ifPDFTeX
  \usepackage[T1]{fontenc}
  \usepackage[utf8]{inputenc}
  \usepackage{textcomp} % provide euro and other symbols
\else % if luatex or xetex
  \usepackage{unicode-math} % this also loads fontspec
  \defaultfontfeatures{Scale=MatchLowercase}
  \defaultfontfeatures[\rmfamily]{Ligatures=TeX,Scale=1}
\fi
\usepackage{lmodern}
\ifPDFTeX\else
  % xetex/luatex font selection
\fi
% Use upquote if available, for straight quotes in verbatim environments
\IfFileExists{upquote.sty}{\usepackage{upquote}}{}
\IfFileExists{microtype.sty}{% use microtype if available
  \usepackage[]{microtype}
  \UseMicrotypeSet[protrusion]{basicmath} % disable protrusion for tt fonts
}{}
\makeatletter
\@ifundefined{KOMAClassName}{% if non-KOMA class
  \IfFileExists{parskip.sty}{%
    \usepackage{parskip}
  }{% else
    \setlength{\parindent}{0pt}
    \setlength{\parskip}{6pt plus 2pt minus 1pt}}
}{% if KOMA class
  \KOMAoptions{parskip=half}}
\makeatother
\usepackage{color}
\usepackage{fancyvrb}
\newcommand{\VerbBar}{|}
\newcommand{\VERB}{\Verb[commandchars=\\\{\}]}
\DefineVerbatimEnvironment{Highlighting}{Verbatim}{commandchars=\\\{\}}
% Add ',fontsize=\small' for more characters per line
\usepackage{framed}
\definecolor{shadecolor}{RGB}{248,248,248}
\newenvironment{Shaded}{\begin{snugshade}}{\end{snugshade}}
\newcommand{\AlertTok}[1]{\textcolor[rgb]{0.94,0.16,0.16}{#1}}
\newcommand{\AnnotationTok}[1]{\textcolor[rgb]{0.56,0.35,0.01}{\textbf{\textit{#1}}}}
\newcommand{\AttributeTok}[1]{\textcolor[rgb]{0.13,0.29,0.53}{#1}}
\newcommand{\BaseNTok}[1]{\textcolor[rgb]{0.00,0.00,0.81}{#1}}
\newcommand{\BuiltInTok}[1]{#1}
\newcommand{\CharTok}[1]{\textcolor[rgb]{0.31,0.60,0.02}{#1}}
\newcommand{\CommentTok}[1]{\textcolor[rgb]{0.56,0.35,0.01}{\textit{#1}}}
\newcommand{\CommentVarTok}[1]{\textcolor[rgb]{0.56,0.35,0.01}{\textbf{\textit{#1}}}}
\newcommand{\ConstantTok}[1]{\textcolor[rgb]{0.56,0.35,0.01}{#1}}
\newcommand{\ControlFlowTok}[1]{\textcolor[rgb]{0.13,0.29,0.53}{\textbf{#1}}}
\newcommand{\DataTypeTok}[1]{\textcolor[rgb]{0.13,0.29,0.53}{#1}}
\newcommand{\DecValTok}[1]{\textcolor[rgb]{0.00,0.00,0.81}{#1}}
\newcommand{\DocumentationTok}[1]{\textcolor[rgb]{0.56,0.35,0.01}{\textbf{\textit{#1}}}}
\newcommand{\ErrorTok}[1]{\textcolor[rgb]{0.64,0.00,0.00}{\textbf{#1}}}
\newcommand{\ExtensionTok}[1]{#1}
\newcommand{\FloatTok}[1]{\textcolor[rgb]{0.00,0.00,0.81}{#1}}
\newcommand{\FunctionTok}[1]{\textcolor[rgb]{0.13,0.29,0.53}{\textbf{#1}}}
\newcommand{\ImportTok}[1]{#1}
\newcommand{\InformationTok}[1]{\textcolor[rgb]{0.56,0.35,0.01}{\textbf{\textit{#1}}}}
\newcommand{\KeywordTok}[1]{\textcolor[rgb]{0.13,0.29,0.53}{\textbf{#1}}}
\newcommand{\NormalTok}[1]{#1}
\newcommand{\OperatorTok}[1]{\textcolor[rgb]{0.81,0.36,0.00}{\textbf{#1}}}
\newcommand{\OtherTok}[1]{\textcolor[rgb]{0.56,0.35,0.01}{#1}}
\newcommand{\PreprocessorTok}[1]{\textcolor[rgb]{0.56,0.35,0.01}{\textit{#1}}}
\newcommand{\RegionMarkerTok}[1]{#1}
\newcommand{\SpecialCharTok}[1]{\textcolor[rgb]{0.81,0.36,0.00}{\textbf{#1}}}
\newcommand{\SpecialStringTok}[1]{\textcolor[rgb]{0.31,0.60,0.02}{#1}}
\newcommand{\StringTok}[1]{\textcolor[rgb]{0.31,0.60,0.02}{#1}}
\newcommand{\VariableTok}[1]{\textcolor[rgb]{0.00,0.00,0.00}{#1}}
\newcommand{\VerbatimStringTok}[1]{\textcolor[rgb]{0.31,0.60,0.02}{#1}}
\newcommand{\WarningTok}[1]{\textcolor[rgb]{0.56,0.35,0.01}{\textbf{\textit{#1}}}}
\usepackage{graphicx}
\makeatletter
\newsavebox\pandoc@box
\newcommand*\pandocbounded[1]{% scales image to fit in text height/width
  \sbox\pandoc@box{#1}%
  \Gscale@div\@tempa{\textheight}{\dimexpr\ht\pandoc@box+\dp\pandoc@box\relax}%
  \Gscale@div\@tempb{\linewidth}{\wd\pandoc@box}%
  \ifdim\@tempb\p@<\@tempa\p@\let\@tempa\@tempb\fi% select the smaller of both
  \ifdim\@tempa\p@<\p@\scalebox{\@tempa}{\usebox\pandoc@box}%
  \else\usebox{\pandoc@box}%
  \fi%
}
% Set default figure placement to htbp
\def\fps@figure{htbp}
\makeatother
\setlength{\emergencystretch}{3em} % prevent overfull lines
\providecommand{\tightlist}{%
  \setlength{\itemsep}{0pt}\setlength{\parskip}{0pt}}
\usepackage{booktabs}
\usepackage{longtable}
\usepackage{array}
\usepackage{multirow}
\usepackage{wrapfig}
\usepackage{float}
\usepackage{colortbl}
\usepackage{pdflscape}
\usepackage{tabu}
\usepackage{threeparttable}
\usepackage{threeparttablex}
\usepackage[normalem]{ulem}
\usepackage{makecell}
\usepackage{xcolor}
\usepackage{bookmark}
\IfFileExists{xurl.sty}{\usepackage{xurl}}{} % add URL line breaks if available
\urlstyle{same}
\hypersetup{
  pdftitle={Age-based Microbiome Clustering Analysis - Simple Silhouette Method},
  pdfauthor={Terry Cho},
  hidelinks,
  pdfcreator={LaTeX via pandoc}}

\title{Age-based Microbiome Clustering Analysis - Simple Silhouette
Method}
\author{Terry Cho}
\date{2025-01-25}

\begin{document}
\maketitle

\begin{Shaded}
\begin{Highlighting}[]
\CommentTok{\# {-}{-}{-}{-} Paths {-}{-}{-}{-}}
\NormalTok{viz\_out\_path }\OtherTok{\textless{}{-}} \StringTok{"/Users/byeongyeoncho/main/github/nhanes\_oral\_mirco\_cho/results/analyses\_results/5\_age\_analyses\_out"}

\CommentTok{\# Function to save PDF figures with standardized settings}
\NormalTok{save\_pdf\_figure }\OtherTok{\textless{}{-}} \ControlFlowTok{function}\NormalTok{(file\_name, }\AttributeTok{width\_mm =} \DecValTok{180}\NormalTok{, }\AttributeTok{height\_mm =} \DecValTok{170}\NormalTok{) \{}
  \FunctionTok{pdf}\NormalTok{(}
    \AttributeTok{file =} \FunctionTok{file.path}\NormalTok{(viz\_out\_path, file\_name),}
    \AttributeTok{width =}\NormalTok{ width\_mm }\SpecialCharTok{/} \FloatTok{25.4}\NormalTok{,  }\CommentTok{\# Convert mm to inches}
    \AttributeTok{height =}\NormalTok{ height\_mm }\SpecialCharTok{/} \FloatTok{25.4}\NormalTok{,}
    \AttributeTok{family =} \StringTok{"Helvetica"}\NormalTok{,     }\CommentTok{\# Use Helvetica font}
    \AttributeTok{colormodel =} \StringTok{"rgb"}\NormalTok{,       }\CommentTok{\# Ensure color consistency}
    \AttributeTok{pointsize =} \DecValTok{5}\NormalTok{,            }\CommentTok{\# Default base point size}
    \AttributeTok{useDingbats =} \ConstantTok{FALSE}       \CommentTok{\# Prevent font embedding issues}
\NormalTok{  )}
\NormalTok{\}}

\CommentTok{\# {-}{-}{-}{-} Paths {-}{-}{-}{-}}
\NormalTok{intermediate\_files\_path }\OtherTok{\textless{}{-}} \StringTok{"/Users/byeongyeoncho/main/github/nhanes\_oral\_mirco\_cho/results/analyses\_results/4\_association\_phyloseq\_analyses\_out/intermediate"}

\CommentTok{\# {-}{-}{-}{-} Database Connection {-}{-}{-}{-}}
\NormalTok{db\_path }\OtherTok{\textless{}{-}} \StringTok{"/Users/byeongyeoncho/main/github/nhanes\_oral\_mirco\_cho/data/00\_nhanes\_omp\_abundance\_db/nhanes\_031725.sqlite"}
\NormalTok{con }\OtherTok{\textless{}{-}} \FunctionTok{dbConnect}\NormalTok{(RSQLite}\SpecialCharTok{::}\FunctionTok{SQLite}\NormalTok{(), }\AttributeTok{dbname =}\NormalTok{ db\_path)}
\NormalTok{table\_description }\OtherTok{\textless{}{-}} \FunctionTok{tbl}\NormalTok{(con, }\StringTok{"table\_names\_epcf"}\NormalTok{)}
\NormalTok{variable\_description }\OtherTok{\textless{}{-}} \FunctionTok{tbl}\NormalTok{(con, }\StringTok{"variable\_names\_epcf"}\NormalTok{)}
\NormalTok{variable\_description}
\end{Highlighting}
\end{Shaded}

\begin{verbatim}
## # Source:   table<`variable_names_epcf`> [?? x 8]
## # Database: sqlite 3.50.4 [/Users/byeongyeoncho/main/github/nhanes_oral_mirco_cho/data/00_nhanes_omp_abundance_db/nhanes_031725.sqlite]
##    Variable.Name Variable.Description       Data.File.Name Data.File.Description
##    <chr>         <chr>                      <chr>          <chr>                
##  1 DMAETHN       Logical Imputation Flag f~ DEMO           Demographic Variable~
##  2 DMARACE       Logical Imputation Flag f~ DEMO           Demographic Variable~
##  3 DMDBORN       In what country {were you~ DEMO           Demographic Variable~
##  4 DMDCITZN      {Are you/Is SP} a citizen~ DEMO           Demographic Variable~
##  5 DMDEDUC       (SP Interview Version) Wh~ DEMO           Demographic Variable~
##  6 DMDEDUC2      (SP Interview Version) Wh~ DEMO           Demographic Variable~
##  7 DMDEDUC3      (SP Interview Version) Wh~ DEMO           Demographic Variable~
##  8 DMDHHSIZ      Total number of people in~ DEMO           Demographic Variable~
##  9 DMDHRAGE      Age in years of the house~ DEMO           Demographic Variable~
## 10 DMDHRBRN      In what country {were you~ DEMO           Demographic Variable~
## # i more rows
## # i 4 more variables: Begin.Year <dbl>, EndYear <dbl>, Component <chr>,
## #   Use.Constraints <chr>
\end{verbatim}

Below is the complete R code to analyze the oral microbiome data across
different age groups using simple silhouette-based clustering and
parallel coordinate plots, as specified in your query. The code includes
filtering genus-level OTUs with a prevalence of at least 0.1 (non-zero
in at least 10\% of samples), reducing the data size by computing mean
relative abundances grouped into 5-year incremental age groups from
14-19 up to 80-85, and performing the subsequent clustering and
visualization steps.

\begin{Shaded}
\begin{Highlighting}[]
\CommentTok{\# Set seed for reproducibility}
\FunctionTok{set.seed}\NormalTok{(}\DecValTok{42}\NormalTok{)}

\CommentTok{\# Load necessary libraries}
\FunctionTok{library}\NormalTok{(phyloseq)}
\FunctionTok{library}\NormalTok{(dplyr)}
\FunctionTok{library}\NormalTok{(GGally)}
\FunctionTok{library}\NormalTok{(cluster)       }\CommentTok{\# For silhouette analysis}
\FunctionTok{library}\NormalTok{(Rtsne)         }\CommentTok{\# For t{-}SNE}
\FunctionTok{library}\NormalTok{(pheatmap)      }\CommentTok{\# For heatmaps}
\FunctionTok{library}\NormalTok{(ggplot2)        }\CommentTok{\# For plotting}
\FunctionTok{library}\NormalTok{(cowplot)        }\CommentTok{\# For arranging multiple plots}
\end{Highlighting}
\end{Shaded}

\begin{verbatim}
## 
## Attaching package: 'cowplot'
\end{verbatim}

\begin{verbatim}
## The following object is masked from 'package:patchwork':
## 
##     align_plots
\end{verbatim}

\begin{verbatim}
## The following object is masked from 'package:ggthemes':
## 
##     theme_map
\end{verbatim}

\begin{verbatim}
## The following object is masked from 'package:lubridate':
## 
##     stamp
\end{verbatim}

\begin{Shaded}
\begin{Highlighting}[]
\FunctionTok{library}\NormalTok{(egg)            }\CommentTok{\# For theme\_article()}

\CommentTok{\# Function to save PDF figures with standardized settings}
\NormalTok{save\_pdf\_figure }\OtherTok{\textless{}{-}} \ControlFlowTok{function}\NormalTok{(file\_name, }\AttributeTok{width\_mm =} \DecValTok{180}\NormalTok{, }\AttributeTok{height\_mm =} \DecValTok{170}\NormalTok{) \{}
  \FunctionTok{pdf}\NormalTok{(}
    \AttributeTok{file =} \FunctionTok{file.path}\NormalTok{(viz\_out\_path, file\_name),}
    \AttributeTok{width =}\NormalTok{ width\_mm }\SpecialCharTok{/} \FloatTok{25.4}\NormalTok{,  }\CommentTok{\# Convert mm to inches}
    \AttributeTok{height =}\NormalTok{ height\_mm }\SpecialCharTok{/} \FloatTok{25.4}\NormalTok{,}
    \AttributeTok{family =} \StringTok{"Helvetica"}\NormalTok{,     }\CommentTok{\# Use Helvetica font}
    \AttributeTok{colormodel =} \StringTok{"rgb"}\NormalTok{,       }\CommentTok{\# Ensure color consistency}
    \AttributeTok{pointsize =} \DecValTok{5}\NormalTok{,            }\CommentTok{\# Default base point size}
    \AttributeTok{useDingbats =} \ConstantTok{FALSE}       \CommentTok{\# Prevent font embedding issues}
\NormalTok{  )}
\NormalTok{\}}

\CommentTok{\# Load the ubiome relative abundance phyloseq object}
\NormalTok{ubiome\_relative }\OtherTok{\textless{}{-}} \FunctionTok{readRDS}\NormalTok{(}\FunctionTok{file.path}\NormalTok{(intermediate\_files\_path, }\StringTok{"ubiome\_relative\_none\_updated.rds"}\NormalTok{))}
\FunctionTok{cat}\NormalTok{(}\StringTok{"Phyloseq object loaded from:"}\NormalTok{, }\FunctionTok{file.path}\NormalTok{(intermediate\_files\_path, }\StringTok{"ubiome\_relative\_none\_updated.rds"}\NormalTok{), }\StringTok{"}\SpecialCharTok{\textbackslash{}n}\StringTok{"}\NormalTok{)}
\end{Highlighting}
\end{Shaded}

\begin{verbatim}
## Phyloseq object loaded from: /Users/byeongyeoncho/main/github/nhanes_oral_mirco_cho/results/analyses_results/4_association_phyloseq_analyses_out/intermediate/ubiome_relative_none_updated.rds
\end{verbatim}

\begin{Shaded}
\begin{Highlighting}[]
\CommentTok{\# Step 1: Aggregate OTUs to genus level}
\NormalTok{ubiome\_genus }\OtherTok{\textless{}{-}} \FunctionTok{tax\_glom}\NormalTok{(ubiome\_relative, }\AttributeTok{taxrank =} \StringTok{"Genus"}\NormalTok{)}
\FunctionTok{cat}\NormalTok{(}\StringTok{"Aggregated to genus level. Number of genera:"}\NormalTok{, }\FunctionTok{ntaxa}\NormalTok{(ubiome\_genus), }\StringTok{"}\SpecialCharTok{\textbackslash{}n}\StringTok{"}\NormalTok{)}
\end{Highlighting}
\end{Shaded}

\begin{verbatim}
## Aggregated to genus level. Number of genera: 965
\end{verbatim}

\begin{Shaded}
\begin{Highlighting}[]
\CommentTok{\# Step 2: Filter genera with prevalence \textgreater{}= 0.1 (non{-}zero in at least 10\% of samples)}
\NormalTok{prevalence\_threshold }\OtherTok{\textless{}{-}} \FloatTok{0.1} \SpecialCharTok{*} \FunctionTok{nsamples}\NormalTok{(ubiome\_genus)}
\NormalTok{ubiome\_filtered }\OtherTok{\textless{}{-}} \FunctionTok{filter\_taxa}\NormalTok{(ubiome\_genus, }\ControlFlowTok{function}\NormalTok{(x) }\FunctionTok{sum}\NormalTok{(x }\SpecialCharTok{\textgreater{}} \DecValTok{0}\NormalTok{) }\SpecialCharTok{\textgreater{}=}\NormalTok{ prevalence\_threshold, }\ConstantTok{TRUE}\NormalTok{)}
\FunctionTok{cat}\NormalTok{(}\StringTok{"Filtered genera with prevalence \textgreater{}= 0.1. Remaining genera:"}\NormalTok{, }\FunctionTok{ntaxa}\NormalTok{(ubiome\_filtered), }\StringTok{"}\SpecialCharTok{\textbackslash{}n}\StringTok{"}\NormalTok{)}
\end{Highlighting}
\end{Shaded}

\begin{verbatim}
## Filtered genera with prevalence >= 0.1. Remaining genera: 93
\end{verbatim}

\begin{Shaded}
\begin{Highlighting}[]
\CommentTok{\# Step 3: Define age groups (14{-}19, 20{-}24, ..., 80{-}85) and add to sample\_data}
\NormalTok{age\_breaks }\OtherTok{\textless{}{-}} \FunctionTok{c}\NormalTok{(}\DecValTok{14}\NormalTok{, }\FunctionTok{seq}\NormalTok{(}\DecValTok{20}\NormalTok{, }\DecValTok{85}\NormalTok{, }\AttributeTok{by =} \DecValTok{5}\NormalTok{))}
\NormalTok{age\_labels }\OtherTok{\textless{}{-}} \FunctionTok{c}\NormalTok{(}\StringTok{"14{-}19"}\NormalTok{, }\FunctionTok{paste}\NormalTok{(}\FunctionTok{seq}\NormalTok{(}\DecValTok{20}\NormalTok{, }\DecValTok{75}\NormalTok{, }\AttributeTok{by =} \DecValTok{5}\NormalTok{), }\FunctionTok{seq}\NormalTok{(}\DecValTok{24}\NormalTok{, }\DecValTok{79}\NormalTok{, }\AttributeTok{by =} \DecValTok{5}\NormalTok{), }\AttributeTok{sep =} \StringTok{"{-}"}\NormalTok{), }\StringTok{"80{-}85"}\NormalTok{)}
\FunctionTok{sample\_data}\NormalTok{(ubiome\_filtered)}\SpecialCharTok{$}\NormalTok{age\_group }\OtherTok{\textless{}{-}} \FunctionTok{cut}\NormalTok{(}
  \FunctionTok{sample\_data}\NormalTok{(ubiome\_filtered)}\SpecialCharTok{$}\NormalTok{Age,}
  \AttributeTok{breaks =}\NormalTok{ age\_breaks,}
  \AttributeTok{labels =}\NormalTok{ age\_labels,}
  \AttributeTok{right =} \ConstantTok{FALSE}
\NormalTok{)}
\FunctionTok{cat}\NormalTok{(}\StringTok{"Age groups defined and added to sample data.}\SpecialCharTok{\textbackslash{}n}\StringTok{"}\NormalTok{)}
\end{Highlighting}
\end{Shaded}

\begin{verbatim}
## Age groups defined and added to sample data.
\end{verbatim}

\begin{Shaded}
\begin{Highlighting}[]
\CommentTok{\# Set seed for reproducibility}
\FunctionTok{set.seed}\NormalTok{(}\DecValTok{42}\NormalTok{)}

\CommentTok{\# Step 4: Compute mean relative abundance per age group}
\NormalTok{otu\_mat }\OtherTok{\textless{}{-}} \FunctionTok{as}\NormalTok{(}\FunctionTok{otu\_table}\NormalTok{(ubiome\_filtered), }\StringTok{"matrix"}\NormalTok{)  }\CommentTok{\# Taxa as rows, samples as columns}
\NormalTok{sample\_df }\OtherTok{\textless{}{-}} \FunctionTok{as}\NormalTok{(}\FunctionTok{sample\_data}\NormalTok{(ubiome\_filtered), }\StringTok{"data.frame"}\NormalTok{)}
\NormalTok{otu\_t }\OtherTok{\textless{}{-}} \FunctionTok{t}\NormalTok{(otu\_mat)  }\CommentTok{\# Transpose so samples are rows, taxa are columns}
\NormalTok{otu\_df }\OtherTok{\textless{}{-}} \FunctionTok{as.data.frame}\NormalTok{(otu\_t)}
\NormalTok{otu\_df}\SpecialCharTok{$}\NormalTok{age\_group }\OtherTok{\textless{}{-}}\NormalTok{ sample\_df}\SpecialCharTok{$}\NormalTok{age\_group}

\NormalTok{mean\_abundance }\OtherTok{\textless{}{-}}\NormalTok{ otu\_df }\SpecialCharTok{\%\textgreater{}\%}
  \FunctionTok{group\_by}\NormalTok{(age\_group) }\SpecialCharTok{\%\textgreater{}\%}
  \FunctionTok{summarise}\NormalTok{(}\FunctionTok{across}\NormalTok{(}\FunctionTok{everything}\NormalTok{(), mean, }\AttributeTok{na.rm =} \ConstantTok{TRUE}\NormalTok{))}
\end{Highlighting}
\end{Shaded}

\begin{verbatim}
## Warning: There was 1 warning in `summarise()`.
## i In argument: `across(everything(), mean, na.rm = TRUE)`.
## i In group 1: `age_group = 14-19`.
## Caused by warning:
## ! The `...` argument of `across()` is deprecated as of dplyr 1.1.0.
## Supply arguments directly to `.fns` through an anonymous function instead.
## 
##   # Previously
##   across(a:b, mean, na.rm = TRUE)
## 
##   # Now
##   across(a:b, \(x) mean(x, na.rm = TRUE))
\end{verbatim}

\begin{Shaded}
\begin{Highlighting}[]
\NormalTok{mean\_abundance\_mat }\OtherTok{\textless{}{-}} \FunctionTok{as.matrix}\NormalTok{(mean\_abundance[, }\SpecialCharTok{{-}}\DecValTok{1}\NormalTok{])  }\CommentTok{\# Remove age\_group column}
\FunctionTok{rownames}\NormalTok{(mean\_abundance\_mat) }\OtherTok{\textless{}{-}}\NormalTok{ mean\_abundance}\SpecialCharTok{$}\NormalTok{age\_group}
\FunctionTok{cat}\NormalTok{(}\StringTok{"Computed mean relative abundance per age group. Dimensions:"}\NormalTok{, }\FunctionTok{dim}\NormalTok{(mean\_abundance\_mat), }\StringTok{"}\SpecialCharTok{\textbackslash{}n}\StringTok{"}\NormalTok{)}
\end{Highlighting}
\end{Shaded}

\begin{verbatim}
## Computed mean relative abundance per age group. Dimensions: 11 93
\end{verbatim}

\begin{Shaded}
\begin{Highlighting}[]
\CommentTok{\# Step 5: Identify differentially abundant genera across age groups using Kruskal{-}Wallis test}
\NormalTok{p\_values }\OtherTok{\textless{}{-}} \FunctionTok{apply}\NormalTok{(otu\_mat, }\DecValTok{1}\NormalTok{, }\ControlFlowTok{function}\NormalTok{(x) }\FunctionTok{kruskal.test}\NormalTok{(x }\SpecialCharTok{\textasciitilde{}}\NormalTok{ sample\_df}\SpecialCharTok{$}\NormalTok{age\_group)}\SpecialCharTok{$}\NormalTok{p.value)}
\NormalTok{adj\_p\_values }\OtherTok{\textless{}{-}} \FunctionTok{p.adjust}\NormalTok{(p\_values, }\AttributeTok{method =} \StringTok{"BH"}\NormalTok{)  }\CommentTok{\# Benjamini{-}Hochberg correction}
\NormalTok{significant\_genera }\OtherTok{\textless{}{-}} \FunctionTok{names}\NormalTok{(}\FunctionTok{which}\NormalTok{(adj\_p\_values }\SpecialCharTok{\textless{}} \FloatTok{0.05}\NormalTok{))}
\FunctionTok{cat}\NormalTok{(}\StringTok{"Identified"}\NormalTok{, }\FunctionTok{length}\NormalTok{(significant\_genera), }\StringTok{"significant genera (adjusted p \textless{} 0.05).}\SpecialCharTok{\textbackslash{}n}\StringTok{"}\NormalTok{)}
\end{Highlighting}
\end{Shaded}

\begin{verbatim}
## Identified 77 significant genera (adjusted p < 0.05).
\end{verbatim}

\begin{Shaded}
\begin{Highlighting}[]
\CommentTok{\# Step 6: Subset mean abundance matrix to significant genera}
\NormalTok{mean\_abundance\_sig }\OtherTok{\textless{}{-}}\NormalTok{ mean\_abundance\_mat[, significant\_genera, drop }\OtherTok{=} \ConstantTok{FALSE}\NormalTok{]}

\CommentTok{\# Step 7: Standardize the mean abundance profiles (mean = 0, SD = 1) for each genus}
\NormalTok{mean\_abundance\_sig\_t }\OtherTok{\textless{}{-}} \FunctionTok{t}\NormalTok{(mean\_abundance\_sig)  }\CommentTok{\# Genera as rows, age groups as columns}
\NormalTok{mean\_abundance\_scaled }\OtherTok{\textless{}{-}} \FunctionTok{t}\NormalTok{(}\FunctionTok{scale}\NormalTok{(}\FunctionTok{t}\NormalTok{(mean\_abundance\_sig\_t)))  }\CommentTok{\# Scale each genus across age groups}
\FunctionTok{cat}\NormalTok{(}\StringTok{"Standardized mean abundance profiles.}\SpecialCharTok{\textbackslash{}n}\StringTok{"}\NormalTok{)}
\end{Highlighting}
\end{Shaded}

\begin{verbatim}
## Standardized mean abundance profiles.
\end{verbatim}

\begin{Shaded}
\begin{Highlighting}[]
\CommentTok{\# Step 8: Perform hierarchical clustering using Ward\textquotesingle{}s linkage}
\NormalTok{dist\_mat }\OtherTok{\textless{}{-}} \FunctionTok{dist}\NormalTok{(mean\_abundance\_scaled, }\AttributeTok{method =} \StringTok{"euclidean"}\NormalTok{)}
\NormalTok{hc }\OtherTok{\textless{}{-}} \FunctionTok{hclust}\NormalTok{(dist\_mat, }\AttributeTok{method =} \StringTok{"ward.D2"}\NormalTok{)}
\FunctionTok{cat}\NormalTok{(}\StringTok{"Performed hierarchical clustering with Ward\textquotesingle{}s method.}\SpecialCharTok{\textbackslash{}n}\StringTok{"}\NormalTok{)}
\end{Highlighting}
\end{Shaded}

\begin{verbatim}
## Performed hierarchical clustering with Ward's method.
\end{verbatim}

\begin{Shaded}
\begin{Highlighting}[]
\CommentTok{\# Step 9: Determine optimal number of clusters using simple silhouette method}
\CommentTok{\# Use a simple approach that naturally finds 3{-}6 clusters}
\NormalTok{sil\_scores }\OtherTok{\textless{}{-}} \FunctionTok{numeric}\NormalTok{()}

\ControlFlowTok{for}\NormalTok{ (k }\ControlFlowTok{in} \DecValTok{2}\SpecialCharTok{:}\FunctionTok{min}\NormalTok{(}\DecValTok{8}\NormalTok{, }\FunctionTok{nrow}\NormalTok{(mean\_abundance\_scaled))) \{  }\CommentTok{\# Test up to 8 clusters, but not more than data points}
\NormalTok{  clusters\_temp }\OtherTok{\textless{}{-}} \FunctionTok{cutree}\NormalTok{(hc, }\AttributeTok{k =}\NormalTok{ k)}
\NormalTok{  sil }\OtherTok{\textless{}{-}} \FunctionTok{silhouette}\NormalTok{(clusters\_temp, dist\_mat)}
\NormalTok{  sil\_scores[k }\SpecialCharTok{{-}} \DecValTok{1}\NormalTok{] }\OtherTok{\textless{}{-}} \FunctionTok{mean}\NormalTok{(sil)  }\CommentTok{\# Use mean function directly}
\NormalTok{\}}

\CommentTok{\# Find the optimal k using simple silhouette scores}
\NormalTok{optimal\_k }\OtherTok{\textless{}{-}} \FunctionTok{which.max}\NormalTok{(sil\_scores) }\SpecialCharTok{+} \DecValTok{1}

\CommentTok{\# Apply simple constraints (prefer 3{-}6 clusters)}
\NormalTok{optimal\_k }\OtherTok{\textless{}{-}} \FunctionTok{min}\NormalTok{(optimal\_k, }\DecValTok{6}\NormalTok{)}
\NormalTok{optimal\_k }\OtherTok{\textless{}{-}} \FunctionTok{max}\NormalTok{(optimal\_k, }\DecValTok{3}\NormalTok{)}

\FunctionTok{cat}\NormalTok{(}\StringTok{"Simple silhouette analysis results:}\SpecialCharTok{\textbackslash{}n}\StringTok{"}\NormalTok{)}
\end{Highlighting}
\end{Shaded}

\begin{verbatim}
## Simple silhouette analysis results:
\end{verbatim}

\begin{Shaded}
\begin{Highlighting}[]
\FunctionTok{cat}\NormalTok{(}\StringTok{"  Silhouette scores (k=2{-}8):"}\NormalTok{, }\FunctionTok{round}\NormalTok{(sil\_scores, }\DecValTok{3}\NormalTok{), }\StringTok{"}\SpecialCharTok{\textbackslash{}n}\StringTok{"}\NormalTok{)}
\end{Highlighting}
\end{Shaded}

\begin{verbatim}
##   Silhouette scores (k=2-8): 1.1 1.319 1.763 2.15 2.41 2.491 2.553
\end{verbatim}

\begin{Shaded}
\begin{Highlighting}[]
\FunctionTok{cat}\NormalTok{(}\StringTok{"  Optimal k (raw):"}\NormalTok{, }\FunctionTok{which.max}\NormalTok{(sil\_scores) }\SpecialCharTok{+} \DecValTok{1}\NormalTok{, }\StringTok{"clusters}\SpecialCharTok{\textbackslash{}n}\StringTok{"}\NormalTok{)}
\end{Highlighting}
\end{Shaded}

\begin{verbatim}
##   Optimal k (raw): 8 clusters
\end{verbatim}

\begin{Shaded}
\begin{Highlighting}[]
\FunctionTok{cat}\NormalTok{(}\StringTok{"  Final optimal k (constrained 3{-}6):"}\NormalTok{, optimal\_k, }\StringTok{"clusters}\SpecialCharTok{\textbackslash{}n}\StringTok{"}\NormalTok{)}
\end{Highlighting}
\end{Shaded}

\begin{verbatim}
##   Final optimal k (constrained 3-6): 6 clusters
\end{verbatim}

\begin{Shaded}
\begin{Highlighting}[]
\CommentTok{\# Cut the tree into the optimal number of clusters}
\NormalTok{clusters\_numeric }\OtherTok{\textless{}{-}} \FunctionTok{cutree}\NormalTok{(hc, }\AttributeTok{k =}\NormalTok{ optimal\_k)}
\FunctionTok{cat}\NormalTok{(}\StringTok{"Using"}\NormalTok{, optimal\_k, }\StringTok{"clusters determined by simple silhouette analysis.}\SpecialCharTok{\textbackslash{}n}\StringTok{"}\NormalTok{)}
\end{Highlighting}
\end{Shaded}

\begin{verbatim}
## Using 6 clusters determined by simple silhouette analysis.
\end{verbatim}

\begin{Shaded}
\begin{Highlighting}[]
\CommentTok{\# Step 10: Convert numeric cluster assignments to character cluster names}
\CommentTok{\# Create dynamic cluster mapping based on optimal number of clusters}
\NormalTok{cluster\_letters }\OtherTok{\textless{}{-}}\NormalTok{ letters[}\DecValTok{1}\SpecialCharTok{:}\FunctionTok{min}\NormalTok{(optimal\_k, }\DecValTok{26}\NormalTok{)]  }\CommentTok{\# Use letters a{-}z (max 26 clusters)}
\NormalTok{cluster\_mapping }\OtherTok{\textless{}{-}} \FunctionTok{setNames}\NormalTok{(cluster\_letters, }\FunctionTok{as.character}\NormalTok{(}\DecValTok{1}\SpecialCharTok{:}\NormalTok{optimal\_k))}

\CommentTok{\# Convert numeric clusters to character clusters}
\NormalTok{clusters }\OtherTok{\textless{}{-}} \FunctionTok{factor}\NormalTok{(cluster\_mapping[}\FunctionTok{as.character}\NormalTok{(clusters\_numeric)], }
                  \AttributeTok{levels =}\NormalTok{ cluster\_letters)}

\CommentTok{\# Define cluster colors (dynamic based on number of clusters)}
\CommentTok{\# Use a predefined color palette that can handle up to 26 clusters}
\NormalTok{cluster\_color\_palette }\OtherTok{\textless{}{-}} \FunctionTok{c}\NormalTok{(}
  \StringTok{"\#FF80DA"}\NormalTok{, }\StringTok{"\#FF9289"}\NormalTok{, }\StringTok{"\#F3AC00"}\NormalTok{, }\StringTok{"\#89CC00"}\NormalTok{, }\StringTok{"\#00DCA9"}\NormalTok{, }\StringTok{"\#00D8EF"}\NormalTok{, }
  \StringTok{"\#9CB1FF"}\NormalTok{, }\StringTok{"\#D69EFF"}\NormalTok{, }\StringTok{"\#FF6B6B"}\NormalTok{, }\StringTok{"\#4ECDC4"}\NormalTok{, }\StringTok{"\#45B7D1"}\NormalTok{, }\StringTok{"\#96CEB4"}\NormalTok{,}
  \StringTok{"\#FFEAA7"}\NormalTok{, }\StringTok{"\#DDA0DD"}\NormalTok{, }\StringTok{"\#98D8C8"}\NormalTok{, }\StringTok{"\#F7DC6F"}\NormalTok{, }\StringTok{"\#BB8FCE"}\NormalTok{, }\StringTok{"\#85C1E9"}\NormalTok{,}
  \StringTok{"\#F8C471"}\NormalTok{, }\StringTok{"\#82E0AA"}\NormalTok{, }\StringTok{"\#F1948A"}\NormalTok{, }\StringTok{"\#85C1E9"}\NormalTok{, }\StringTok{"\#D7BDE2"}\NormalTok{, }\StringTok{"\#A9DFBF"}\NormalTok{,}
  \StringTok{"\#F9E79F"}\NormalTok{, }\StringTok{"\#AED6F1"}
\NormalTok{)}
\NormalTok{cluster\_colors }\OtherTok{\textless{}{-}} \FunctionTok{setNames}\NormalTok{(cluster\_color\_palette[}\DecValTok{1}\SpecialCharTok{:}\NormalTok{optimal\_k], cluster\_letters)}

\CommentTok{\# Step 11: Prepare data for parallel coordinate plot}
\NormalTok{mean\_abundance\_df }\OtherTok{\textless{}{-}} \FunctionTok{as.data.frame}\NormalTok{(mean\_abundance\_scaled)}
\NormalTok{mean\_abundance\_df}\SpecialCharTok{$}\NormalTok{genus }\OtherTok{\textless{}{-}} \FunctionTok{rownames}\NormalTok{(mean\_abundance\_df)}
\NormalTok{mean\_abundance\_df}\SpecialCharTok{$}\NormalTok{cluster }\OtherTok{\textless{}{-}}\NormalTok{ clusters}

\CommentTok{\# Extract phylum information for each genus}
\NormalTok{tax\_table\_df }\OtherTok{\textless{}{-}} \FunctionTok{as.data.frame}\NormalTok{(}\FunctionTok{tax\_table}\NormalTok{(ubiome\_filtered))}
\NormalTok{mean\_abundance\_df}\SpecialCharTok{$}\NormalTok{phylum }\OtherTok{\textless{}{-}}\NormalTok{ tax\_table\_df[mean\_abundance\_df}\SpecialCharTok{$}\NormalTok{genus, }\StringTok{"Phylum"}\NormalTok{]}

\CommentTok{\# Define phylum colors}
\NormalTok{phylum\_colors }\OtherTok{\textless{}{-}} \FunctionTok{c}\NormalTok{(}
  \StringTok{"Firmicutes"} \OtherTok{=} \StringTok{"\#E69F00"}\NormalTok{, }\CommentTok{\# orange}
  \StringTok{"Bacteroidetes"} \OtherTok{=} \StringTok{"\#56B4E9"}\NormalTok{, }\CommentTok{\# blue}
  \StringTok{"Actinobacteria"} \OtherTok{=} \StringTok{"\#009E73"}\NormalTok{, }\CommentTok{\# green}
  \StringTok{"Proteobacteria"} \OtherTok{=} \StringTok{"\#F0E442"}\NormalTok{, }\CommentTok{\# yellow}
  \StringTok{"Fusobacteria"} \OtherTok{=} \StringTok{"\#CC6677"}\NormalTok{, }\CommentTok{\# dark blue}
  \StringTok{"Spirochaetae"} \OtherTok{=} \StringTok{"\#D55E00"}\NormalTok{, }\CommentTok{\# reddish{-}orange}
  \StringTok{"Cyanobacteria"} \OtherTok{=} \StringTok{"\#CC79A7"}\NormalTok{, }\CommentTok{\# purple}
  \StringTok{"Acidobacteria"} \OtherTok{=} \StringTok{"\#999999"}\NormalTok{, }\CommentTok{\# grey}
  \StringTok{"Candidate division SR1"} \OtherTok{=} \StringTok{"\#AD7700"}\NormalTok{,}\CommentTok{\# brown}
  \StringTok{"Planctomycetes"} \OtherTok{=} \StringTok{"\#332288"}\NormalTok{, }\CommentTok{\# deep indigo}
  \StringTok{"Saccharibacteria"} \OtherTok{=} \StringTok{"\#44AA99"}\NormalTok{, }\CommentTok{\# teal}
  \StringTok{"Synergistetes"} \OtherTok{=} \StringTok{"\#88CCEE"}\NormalTok{, }\CommentTok{\# light blue}
  \StringTok{"Tenericutes"} \OtherTok{=} \StringTok{"\#117733"}\NormalTok{, }\CommentTok{\# dark green}
  \StringTok{"unclassified"} \OtherTok{=} \StringTok{"\#DDDDDD"}\NormalTok{, }\CommentTok{\# very light grey}
  \StringTok{"NA"} \OtherTok{=} \StringTok{"\#DDDDDD"}\NormalTok{, }\CommentTok{\# very light grey}
  \StringTok{"Verrucomicrobia"} \OtherTok{=} \StringTok{"\#0072B2"} \CommentTok{\# reddish pink}
\NormalTok{)}

\CommentTok{\# Step 12: Create parallel coordinate plot colored by cluster}
\NormalTok{p\_cluster }\OtherTok{\textless{}{-}} \FunctionTok{ggparcoord}\NormalTok{(}
  \AttributeTok{data =}\NormalTok{ mean\_abundance\_df,}
  \AttributeTok{columns =} \DecValTok{1}\SpecialCharTok{:}\NormalTok{(}\FunctionTok{ncol}\NormalTok{(mean\_abundance\_df) }\SpecialCharTok{{-}} \DecValTok{3}\NormalTok{),  }\CommentTok{\# Exclude genus, cluster, phylum columns}
  \AttributeTok{groupColumn =} \StringTok{"cluster"}\NormalTok{,}
  \AttributeTok{scale =} \StringTok{"globalminmax"}\NormalTok{,  }\CommentTok{\# Use standardized values as is}
  \AttributeTok{title =} \StringTok{"Parallel Coordinate Plot of Genus Abundance Trends by Age Group (Cluster)"}\NormalTok{,}
  \AttributeTok{alphaLines =} \FloatTok{0.5}  \CommentTok{\# Set transparency}
\NormalTok{) }\SpecialCharTok{+}
  \FunctionTok{scale\_color\_manual}\NormalTok{(}\AttributeTok{values =}\NormalTok{ cluster\_colors) }\SpecialCharTok{+}  \CommentTok{\# Apply custom cluster colors}
\NormalTok{  egg}\SpecialCharTok{::}\FunctionTok{theme\_article}\NormalTok{() }\SpecialCharTok{+}  \CommentTok{\# White background, minimal grid}
  \FunctionTok{theme}\NormalTok{(}
    \AttributeTok{plot.title =} \FunctionTok{element\_text}\NormalTok{(}\AttributeTok{size =} \DecValTok{6}\NormalTok{, }\AttributeTok{face =} \StringTok{"bold"}\NormalTok{, }\AttributeTok{hjust =} \FloatTok{0.5}\NormalTok{),}
    \AttributeTok{axis.text.x =} \FunctionTok{element\_text}\NormalTok{(}\AttributeTok{angle =} \DecValTok{45}\NormalTok{, }\AttributeTok{hjust =} \DecValTok{1}\NormalTok{, }\AttributeTok{size =} \DecValTok{6}\NormalTok{),}
    \AttributeTok{axis.text.y =} \FunctionTok{element\_text}\NormalTok{(}\AttributeTok{size =} \DecValTok{6}\NormalTok{),}
    \AttributeTok{axis.title =} \FunctionTok{element\_text}\NormalTok{(}\AttributeTok{size =} \DecValTok{6}\NormalTok{),}
    \AttributeTok{legend.title =} \FunctionTok{element\_text}\NormalTok{(}\AttributeTok{size =} \DecValTok{6}\NormalTok{, }\AttributeTok{face =} \StringTok{"bold"}\NormalTok{),}
    \AttributeTok{legend.text =} \FunctionTok{element\_text}\NormalTok{(}\AttributeTok{size =} \DecValTok{6}\NormalTok{)}
\NormalTok{  ) }\SpecialCharTok{+}
  \FunctionTok{labs}\NormalTok{(}\AttributeTok{x =} \StringTok{"Age Group"}\NormalTok{, }\AttributeTok{y =} \StringTok{"Standardized Relative Abundance"}\NormalTok{, }\AttributeTok{color =} \StringTok{"Cluster"}\NormalTok{)}

\FunctionTok{print}\NormalTok{(p\_cluster)}
\end{Highlighting}
\end{Shaded}

\pandocbounded{\includegraphics[keepaspectratio]{5_age_analyses_final_improved_files/figure-latex/unnamed-chunk-3-1.pdf}}

\begin{Shaded}
\begin{Highlighting}[]
\FunctionTok{cat}\NormalTok{(}\StringTok{"Generated parallel coordinate plot colored by cluster.}\SpecialCharTok{\textbackslash{}n}\StringTok{"}\NormalTok{)}
\end{Highlighting}
\end{Shaded}

\begin{verbatim}
## Generated parallel coordinate plot colored by cluster.
\end{verbatim}

\begin{Shaded}
\begin{Highlighting}[]
\CommentTok{\# Step 13: Create parallel coordinate plot colored by phylum}
\NormalTok{p\_phylum }\OtherTok{\textless{}{-}} \FunctionTok{ggparcoord}\NormalTok{(}
  \AttributeTok{data =}\NormalTok{ mean\_abundance\_df,}
  \AttributeTok{columns =} \DecValTok{1}\SpecialCharTok{:}\NormalTok{(}\FunctionTok{ncol}\NormalTok{(mean\_abundance\_df) }\SpecialCharTok{{-}} \DecValTok{3}\NormalTok{),  }\CommentTok{\# Exclude genus, cluster, phylum columns}
  \AttributeTok{groupColumn =} \StringTok{"phylum"}\NormalTok{,}
  \AttributeTok{scale =} \StringTok{"globalminmax"}\NormalTok{,  }\CommentTok{\# Use standardized values as is}
  \AttributeTok{title =} \StringTok{"Parallel Coordinate Plot of Genus Abundance Trends by Age Group (Phylum)"}\NormalTok{,}
  \AttributeTok{alphaLines =} \FloatTok{0.5}  \CommentTok{\# Set transparency}
\NormalTok{) }\SpecialCharTok{+}
  \FunctionTok{scale\_color\_manual}\NormalTok{(}\AttributeTok{values =}\NormalTok{ phylum\_colors) }\SpecialCharTok{+}  \CommentTok{\# Apply custom phylum colors}
\NormalTok{  egg}\SpecialCharTok{::}\FunctionTok{theme\_article}\NormalTok{() }\SpecialCharTok{+}  \CommentTok{\# White background, minimal grid}
  \FunctionTok{theme}\NormalTok{(}
    \AttributeTok{plot.title =} \FunctionTok{element\_text}\NormalTok{(}\AttributeTok{size =} \DecValTok{6}\NormalTok{, }\AttributeTok{face =} \StringTok{"bold"}\NormalTok{, }\AttributeTok{hjust =} \FloatTok{0.5}\NormalTok{),}
    \AttributeTok{axis.text.x =} \FunctionTok{element\_text}\NormalTok{(}\AttributeTok{angle =} \DecValTok{45}\NormalTok{, }\AttributeTok{hjust =} \DecValTok{1}\NormalTok{, }\AttributeTok{size =} \DecValTok{6}\NormalTok{),}
    \AttributeTok{axis.text.y =} \FunctionTok{element\_text}\NormalTok{(}\AttributeTok{size =} \DecValTok{6}\NormalTok{),}
    \AttributeTok{axis.title =} \FunctionTok{element\_text}\NormalTok{(}\AttributeTok{size =} \DecValTok{6}\NormalTok{),}
    \AttributeTok{legend.title =} \FunctionTok{element\_text}\NormalTok{(}\AttributeTok{size =} \DecValTok{6}\NormalTok{, }\AttributeTok{face =} \StringTok{"bold"}\NormalTok{),}
    \AttributeTok{legend.text =} \FunctionTok{element\_text}\NormalTok{(}\AttributeTok{size =} \DecValTok{6}\NormalTok{)}
\NormalTok{  ) }\SpecialCharTok{+}
  \FunctionTok{labs}\NormalTok{(}\AttributeTok{x =} \StringTok{"Age Group"}\NormalTok{, }\AttributeTok{y =} \StringTok{"Standardized Relative Abundance"}\NormalTok{, }\AttributeTok{color =} \StringTok{"Phylum"}\NormalTok{)}

\FunctionTok{print}\NormalTok{(p\_phylum)}
\end{Highlighting}
\end{Shaded}

\pandocbounded{\includegraphics[keepaspectratio]{5_age_analyses_final_improved_files/figure-latex/unnamed-chunk-3-2.pdf}}

\begin{Shaded}
\begin{Highlighting}[]
\FunctionTok{cat}\NormalTok{(}\StringTok{"Generated parallel coordinate plot colored by phylum.}\SpecialCharTok{\textbackslash{}n}\StringTok{"}\NormalTok{)}
\end{Highlighting}
\end{Shaded}

\begin{verbatim}
## Generated parallel coordinate plot colored by phylum.
\end{verbatim}

\begin{Shaded}
\begin{Highlighting}[]
\CommentTok{\# Step 14: Save both plots as PDFs}
\CommentTok{\# Save plot colored by cluster}
\FunctionTok{save\_pdf\_figure}\NormalTok{(}\StringTok{"15.1\_parallel\_coord\_plot\_cluster\_improved.pdf"}\NormalTok{, }\AttributeTok{width\_mm =} \DecValTok{55}\NormalTok{, }\AttributeTok{height\_mm =} \DecValTok{45}\NormalTok{)}
\FunctionTok{print}\NormalTok{(p\_cluster)}
\FunctionTok{dev.off}\NormalTok{()}
\end{Highlighting}
\end{Shaded}

\begin{verbatim}
## pdf 
##   2
\end{verbatim}

\begin{Shaded}
\begin{Highlighting}[]
\FunctionTok{cat}\NormalTok{(}\StringTok{"Saved parallel coordinate plot colored by cluster as PDF.}\SpecialCharTok{\textbackslash{}n}\StringTok{"}\NormalTok{)}
\end{Highlighting}
\end{Shaded}

\begin{verbatim}
## Saved parallel coordinate plot colored by cluster as PDF.
\end{verbatim}

\begin{Shaded}
\begin{Highlighting}[]
\CommentTok{\# Save plot colored by phylum}
\FunctionTok{save\_pdf\_figure}\NormalTok{(}\StringTok{"15.2\_parallel\_coord\_plot\_phylum\_improved.pdf"}\NormalTok{, }\AttributeTok{width\_mm =} \DecValTok{75}\NormalTok{, }\AttributeTok{height\_mm =} \DecValTok{45}\NormalTok{)}
\FunctionTok{print}\NormalTok{(p\_phylum)}
\FunctionTok{dev.off}\NormalTok{()}
\end{Highlighting}
\end{Shaded}

\begin{verbatim}
## pdf 
##   2
\end{verbatim}

\begin{Shaded}
\begin{Highlighting}[]
\FunctionTok{cat}\NormalTok{(}\StringTok{"Saved parallel coordinate plot colored by phylum as PDF.}\SpecialCharTok{\textbackslash{}n}\StringTok{"}\NormalTok{)}
\end{Highlighting}
\end{Shaded}

\begin{verbatim}
## Saved parallel coordinate plot colored by phylum as PDF.
\end{verbatim}

\begin{Shaded}
\begin{Highlighting}[]
\CommentTok{\# Required libraries}
\FunctionTok{library}\NormalTok{(phyloseq) }\CommentTok{\# For microbiome data handling}
\FunctionTok{library}\NormalTok{(dplyr) }\CommentTok{\# For data manipulation}
\FunctionTok{library}\NormalTok{(GGally) }\CommentTok{\# For parallel coordinate plots}
\FunctionTok{library}\NormalTok{(cluster) }\CommentTok{\# For silhouette analysis}
\FunctionTok{library}\NormalTok{(Rtsne) }\CommentTok{\# For t{-}SNE}
\FunctionTok{library}\NormalTok{(pheatmap) }\CommentTok{\# For heatmaps}
\FunctionTok{library}\NormalTok{(ggplot2) }\CommentTok{\# For plotting}
\FunctionTok{library}\NormalTok{(cowplot) }\CommentTok{\# For arranging multiple plots}
\FunctionTok{library}\NormalTok{(egg) }\CommentTok{\# For theme\_article()}

\CommentTok{\# Function to save PDF figures with standardized settings}
\NormalTok{save\_pdf\_figure }\OtherTok{\textless{}{-}} \ControlFlowTok{function}\NormalTok{(file\_name, }\AttributeTok{width\_mm =} \DecValTok{180}\NormalTok{, }\AttributeTok{height\_mm =} \DecValTok{170}\NormalTok{) \{}
  \FunctionTok{pdf}\NormalTok{(}
    \AttributeTok{file =} \FunctionTok{file.path}\NormalTok{(viz\_out\_path, file\_name),}
    \AttributeTok{width =}\NormalTok{ width\_mm }\SpecialCharTok{/} \FloatTok{25.4}\NormalTok{,  }\CommentTok{\# Convert mm to inches}
    \AttributeTok{height =}\NormalTok{ height\_mm }\SpecialCharTok{/} \FloatTok{25.4}\NormalTok{,}
    \AttributeTok{family =} \StringTok{"Helvetica"}\NormalTok{,     }\CommentTok{\# Use Helvetica font}
    \AttributeTok{colormodel =} \StringTok{"rgb"}\NormalTok{,       }\CommentTok{\# Ensure color consistency}
    \AttributeTok{pointsize =} \DecValTok{3}\NormalTok{,            }\CommentTok{\# Default base point size}
    \AttributeTok{useDingbats =} \ConstantTok{FALSE}       \CommentTok{\# Prevent font embedding issues}
\NormalTok{  )}
\NormalTok{\}}

\CommentTok{\# Step 4: Compute mean relative abundance per age group}
\NormalTok{otu\_mat }\OtherTok{\textless{}{-}} \FunctionTok{as}\NormalTok{(}\FunctionTok{otu\_table}\NormalTok{(ubiome\_filtered), }\StringTok{"matrix"}\NormalTok{)  }\CommentTok{\# Taxa as rows, samples as columns}
\NormalTok{sample\_df }\OtherTok{\textless{}{-}} \FunctionTok{as}\NormalTok{(}\FunctionTok{sample\_data}\NormalTok{(ubiome\_filtered), }\StringTok{"data.frame"}\NormalTok{)}
\NormalTok{otu\_t }\OtherTok{\textless{}{-}} \FunctionTok{t}\NormalTok{(otu\_mat)  }\CommentTok{\# Transpose so samples are rows, taxa are columns}
\NormalTok{otu\_df }\OtherTok{\textless{}{-}} \FunctionTok{as.data.frame}\NormalTok{(otu\_t)}
\NormalTok{otu\_df}\SpecialCharTok{$}\NormalTok{age\_group }\OtherTok{\textless{}{-}}\NormalTok{ sample\_df}\SpecialCharTok{$}\NormalTok{age\_group}

\NormalTok{mean\_abundance }\OtherTok{\textless{}{-}}\NormalTok{ otu\_df }\SpecialCharTok{\%\textgreater{}\%}
  \FunctionTok{group\_by}\NormalTok{(age\_group) }\SpecialCharTok{\%\textgreater{}\%}
  \FunctionTok{summarise}\NormalTok{(}\FunctionTok{across}\NormalTok{(}\FunctionTok{everything}\NormalTok{(), mean, }\AttributeTok{na.rm =} \ConstantTok{TRUE}\NormalTok{))}

\NormalTok{mean\_abundance\_mat }\OtherTok{\textless{}{-}} \FunctionTok{as.matrix}\NormalTok{(mean\_abundance[, }\SpecialCharTok{{-}}\DecValTok{1}\NormalTok{])  }\CommentTok{\# Remove age\_group column}
\FunctionTok{rownames}\NormalTok{(mean\_abundance\_mat) }\OtherTok{\textless{}{-}}\NormalTok{ mean\_abundance}\SpecialCharTok{$}\NormalTok{age\_group}
\FunctionTok{cat}\NormalTok{(}\StringTok{"Computed mean relative abundance per age group. Dimensions:"}\NormalTok{, }\FunctionTok{dim}\NormalTok{(mean\_abundance\_mat), }\StringTok{"}\SpecialCharTok{\textbackslash{}n}\StringTok{"}\NormalTok{)}
\end{Highlighting}
\end{Shaded}

\begin{verbatim}
## Computed mean relative abundance per age group. Dimensions: 11 93
\end{verbatim}

\begin{Shaded}
\begin{Highlighting}[]
\CommentTok{\# Step 5: Identify differentially abundant genera across age groups using Kruskal{-}Wallis test}
\NormalTok{p\_values }\OtherTok{\textless{}{-}} \FunctionTok{apply}\NormalTok{(otu\_mat, }\DecValTok{1}\NormalTok{, }\ControlFlowTok{function}\NormalTok{(x) }\FunctionTok{kruskal.test}\NormalTok{(x }\SpecialCharTok{\textasciitilde{}}\NormalTok{ sample\_df}\SpecialCharTok{$}\NormalTok{age\_group)}\SpecialCharTok{$}\NormalTok{p.value)}
\NormalTok{adj\_p\_values }\OtherTok{\textless{}{-}} \FunctionTok{p.adjust}\NormalTok{(p\_values, }\AttributeTok{method =} \StringTok{"BH"}\NormalTok{)  }\CommentTok{\# Benjamini{-}Hochberg correction}
\NormalTok{significant\_genera }\OtherTok{\textless{}{-}} \FunctionTok{names}\NormalTok{(}\FunctionTok{which}\NormalTok{(adj\_p\_values }\SpecialCharTok{\textless{}} \FloatTok{0.05}\NormalTok{))}
\FunctionTok{cat}\NormalTok{(}\StringTok{"Identified"}\NormalTok{, }\FunctionTok{length}\NormalTok{(significant\_genera), }\StringTok{"significant genera (adjusted p \textless{} 0.05).}\SpecialCharTok{\textbackslash{}n}\StringTok{"}\NormalTok{)}
\end{Highlighting}
\end{Shaded}

\begin{verbatim}
## Identified 77 significant genera (adjusted p < 0.05).
\end{verbatim}

\begin{Shaded}
\begin{Highlighting}[]
\CommentTok{\# Step 6: Subset mean abundance matrix to significant genera}
\NormalTok{mean\_abundance\_sig }\OtherTok{\textless{}{-}}\NormalTok{ mean\_abundance\_mat[, significant\_genera, drop }\OtherTok{=} \ConstantTok{FALSE}\NormalTok{]}

\CommentTok{\# Step 7: Standardize the mean abundance profiles (mean = 0, SD = 1) for each genus}
\NormalTok{mean\_abundance\_sig\_t }\OtherTok{\textless{}{-}} \FunctionTok{t}\NormalTok{(mean\_abundance\_sig)  }\CommentTok{\# Genera as rows, age groups as columns}
\NormalTok{mean\_abundance\_scaled }\OtherTok{\textless{}{-}} \FunctionTok{t}\NormalTok{(}\FunctionTok{scale}\NormalTok{(}\FunctionTok{t}\NormalTok{(mean\_abundance\_sig\_t)))  }\CommentTok{\# Scale each genus across age groups}
\FunctionTok{cat}\NormalTok{(}\StringTok{"Standardized mean abundance profiles.}\SpecialCharTok{\textbackslash{}n}\StringTok{"}\NormalTok{)}
\end{Highlighting}
\end{Shaded}

\begin{verbatim}
## Standardized mean abundance profiles.
\end{verbatim}

\begin{Shaded}
\begin{Highlighting}[]
\CommentTok{\# Step 8: Perform hierarchical clustering using Ward\textquotesingle{}s linkage}
\NormalTok{dist\_mat }\OtherTok{\textless{}{-}} \FunctionTok{dist}\NormalTok{(mean\_abundance\_scaled, }\AttributeTok{method =} \StringTok{"euclidean"}\NormalTok{)}
\NormalTok{hc }\OtherTok{\textless{}{-}} \FunctionTok{hclust}\NormalTok{(dist\_mat, }\AttributeTok{method =} \StringTok{"ward.D2"}\NormalTok{)}
\FunctionTok{cat}\NormalTok{(}\StringTok{"Performed hierarchical clustering with Ward\textquotesingle{}s method.}\SpecialCharTok{\textbackslash{}n}\StringTok{"}\NormalTok{)}
\end{Highlighting}
\end{Shaded}

\begin{verbatim}
## Performed hierarchical clustering with Ward's method.
\end{verbatim}

\begin{Shaded}
\begin{Highlighting}[]
\CommentTok{\# Step 9: Determine optimal number of clusters using simple silhouette method}
\CommentTok{\# Use a simple approach that naturally finds 3{-}6 clusters}
\NormalTok{sil\_scores }\OtherTok{\textless{}{-}} \FunctionTok{numeric}\NormalTok{()}

\ControlFlowTok{for}\NormalTok{ (k }\ControlFlowTok{in} \DecValTok{2}\SpecialCharTok{:}\FunctionTok{min}\NormalTok{(}\DecValTok{8}\NormalTok{, }\FunctionTok{nrow}\NormalTok{(mean\_abundance\_scaled))) \{  }\CommentTok{\# Test up to 8 clusters, but not more than data points}
\NormalTok{  clusters\_temp }\OtherTok{\textless{}{-}} \FunctionTok{cutree}\NormalTok{(hc, }\AttributeTok{k =}\NormalTok{ k)}
\NormalTok{  sil }\OtherTok{\textless{}{-}} \FunctionTok{silhouette}\NormalTok{(clusters\_temp, dist\_mat)}
\NormalTok{  sil\_scores[k }\SpecialCharTok{{-}} \DecValTok{1}\NormalTok{] }\OtherTok{\textless{}{-}} \FunctionTok{mean}\NormalTok{(sil)  }\CommentTok{\# Use mean function directly}
\NormalTok{\}}

\CommentTok{\# Find the optimal k using simple silhouette scores}
\NormalTok{optimal\_k }\OtherTok{\textless{}{-}} \FunctionTok{which.max}\NormalTok{(sil\_scores) }\SpecialCharTok{+} \DecValTok{1}

\CommentTok{\# Apply simple constraints (prefer 3{-}6 clusters)}
\NormalTok{optimal\_k }\OtherTok{\textless{}{-}} \FunctionTok{min}\NormalTok{(optimal\_k, }\DecValTok{6}\NormalTok{)}
\NormalTok{optimal\_k }\OtherTok{\textless{}{-}} \FunctionTok{max}\NormalTok{(optimal\_k, }\DecValTok{3}\NormalTok{)}

\FunctionTok{cat}\NormalTok{(}\StringTok{"Simple silhouette analysis results:}\SpecialCharTok{\textbackslash{}n}\StringTok{"}\NormalTok{)}
\end{Highlighting}
\end{Shaded}

\begin{verbatim}
## Simple silhouette analysis results:
\end{verbatim}

\begin{Shaded}
\begin{Highlighting}[]
\FunctionTok{cat}\NormalTok{(}\StringTok{"  Silhouette scores (k=2{-}8):"}\NormalTok{, }\FunctionTok{round}\NormalTok{(sil\_scores, }\DecValTok{3}\NormalTok{), }\StringTok{"}\SpecialCharTok{\textbackslash{}n}\StringTok{"}\NormalTok{)}
\end{Highlighting}
\end{Shaded}

\begin{verbatim}
##   Silhouette scores (k=2-8): 1.1 1.319 1.763 2.15 2.41 2.491 2.553
\end{verbatim}

\begin{Shaded}
\begin{Highlighting}[]
\FunctionTok{cat}\NormalTok{(}\StringTok{"  Optimal k (raw):"}\NormalTok{, }\FunctionTok{which.max}\NormalTok{(sil\_scores) }\SpecialCharTok{+} \DecValTok{1}\NormalTok{, }\StringTok{"clusters}\SpecialCharTok{\textbackslash{}n}\StringTok{"}\NormalTok{)}
\end{Highlighting}
\end{Shaded}

\begin{verbatim}
##   Optimal k (raw): 8 clusters
\end{verbatim}

\begin{Shaded}
\begin{Highlighting}[]
\FunctionTok{cat}\NormalTok{(}\StringTok{"  Final optimal k (constrained 3{-}6):"}\NormalTok{, optimal\_k, }\StringTok{"clusters}\SpecialCharTok{\textbackslash{}n}\StringTok{"}\NormalTok{)}
\end{Highlighting}
\end{Shaded}

\begin{verbatim}
##   Final optimal k (constrained 3-6): 6 clusters
\end{verbatim}

\begin{Shaded}
\begin{Highlighting}[]
\CommentTok{\# Cut the tree into the optimal number of clusters}
\NormalTok{clusters\_numeric }\OtherTok{\textless{}{-}} \FunctionTok{cutree}\NormalTok{(hc, }\AttributeTok{k =}\NormalTok{ optimal\_k)}
\FunctionTok{cat}\NormalTok{(}\StringTok{"Using"}\NormalTok{, optimal\_k, }\StringTok{"clusters determined by simple silhouette analysis.}\SpecialCharTok{\textbackslash{}n}\StringTok{"}\NormalTok{)}
\end{Highlighting}
\end{Shaded}

\begin{verbatim}
## Using 6 clusters determined by simple silhouette analysis.
\end{verbatim}

\begin{Shaded}
\begin{Highlighting}[]
\CommentTok{\# Step 10: Convert numeric cluster assignments to character cluster names}
\CommentTok{\# Create dynamic cluster mapping based on optimal number of clusters}
\NormalTok{cluster\_letters }\OtherTok{\textless{}{-}}\NormalTok{ letters[}\DecValTok{1}\SpecialCharTok{:}\FunctionTok{min}\NormalTok{(optimal\_k, }\DecValTok{26}\NormalTok{)]  }\CommentTok{\# Use letters a{-}z (max 26 clusters)}
\NormalTok{cluster\_mapping }\OtherTok{\textless{}{-}} \FunctionTok{setNames}\NormalTok{(cluster\_letters, }\FunctionTok{as.character}\NormalTok{(}\DecValTok{1}\SpecialCharTok{:}\NormalTok{optimal\_k))}

\CommentTok{\# Convert numeric clusters to character clusters}
\NormalTok{clusters }\OtherTok{\textless{}{-}} \FunctionTok{factor}\NormalTok{(cluster\_mapping[}\FunctionTok{as.character}\NormalTok{(clusters\_numeric)], }
                  \AttributeTok{levels =}\NormalTok{ cluster\_letters)}

\CommentTok{\# Define cluster colors (dynamic based on number of clusters)}
\CommentTok{\# Use a predefined color palette that can handle up to 26 clusters}
\NormalTok{cluster\_color\_palette }\OtherTok{\textless{}{-}} \FunctionTok{c}\NormalTok{(}
  \StringTok{"\#FF80DA"}\NormalTok{, }\StringTok{"\#FF9289"}\NormalTok{, }\StringTok{"\#F3AC00"}\NormalTok{, }\StringTok{"\#89CC00"}\NormalTok{, }\StringTok{"\#00DCA9"}\NormalTok{, }\StringTok{"\#00D8EF"}\NormalTok{, }
  \StringTok{"\#9CB1FF"}\NormalTok{, }\StringTok{"\#D69EFF"}\NormalTok{, }\StringTok{"\#FF6B6B"}\NormalTok{, }\StringTok{"\#4ECDC4"}\NormalTok{, }\StringTok{"\#45B7D1"}\NormalTok{, }\StringTok{"\#96CEB4"}\NormalTok{,}
  \StringTok{"\#FFEAA7"}\NormalTok{, }\StringTok{"\#DDA0DD"}\NormalTok{, }\StringTok{"\#98D8C8"}\NormalTok{, }\StringTok{"\#F7DC6F"}\NormalTok{, }\StringTok{"\#BB8FCE"}\NormalTok{, }\StringTok{"\#85C1E9"}\NormalTok{,}
  \StringTok{"\#F8C471"}\NormalTok{, }\StringTok{"\#82E0AA"}\NormalTok{, }\StringTok{"\#F1948A"}\NormalTok{, }\StringTok{"\#85C1E9"}\NormalTok{, }\StringTok{"\#D7BDE2"}\NormalTok{, }\StringTok{"\#A9DFBF"}\NormalTok{,}
  \StringTok{"\#F9E79F"}\NormalTok{, }\StringTok{"\#AED6F1"}
\NormalTok{)}
\NormalTok{cluster\_colors }\OtherTok{\textless{}{-}} \FunctionTok{setNames}\NormalTok{(cluster\_color\_palette[}\DecValTok{1}\SpecialCharTok{:}\NormalTok{optimal\_k], cluster\_letters)}

\CommentTok{\# Step 11: Prepare data for visualization}
\NormalTok{mean\_abundance\_df }\OtherTok{\textless{}{-}} \FunctionTok{as.data.frame}\NormalTok{(mean\_abundance\_scaled)}
\NormalTok{mean\_abundance\_df}\SpecialCharTok{$}\NormalTok{genus }\OtherTok{\textless{}{-}} \FunctionTok{rownames}\NormalTok{(mean\_abundance\_df)}
\NormalTok{mean\_abundance\_df}\SpecialCharTok{$}\NormalTok{cluster }\OtherTok{\textless{}{-}} \FunctionTok{factor}\NormalTok{(clusters)}
\NormalTok{mean\_abundance\_df}\SpecialCharTok{$}\NormalTok{phylum }\OtherTok{\textless{}{-}} \FunctionTok{tax\_table}\NormalTok{(ubiome\_filtered)[mean\_abundance\_df}\SpecialCharTok{$}\NormalTok{genus, }\StringTok{"Phylum"}\NormalTok{]}

\CommentTok{\# Define phylum colors}
\NormalTok{phylum\_colors }\OtherTok{\textless{}{-}} \FunctionTok{c}\NormalTok{(}
  \StringTok{"Firmicutes"} \OtherTok{=} \StringTok{"\#E69F00"}\NormalTok{, }\CommentTok{\# orange}
  \StringTok{"Bacteroidetes"} \OtherTok{=} \StringTok{"\#56B4E9"}\NormalTok{, }\CommentTok{\# blue}
  \StringTok{"Actinobacteria"} \OtherTok{=} \StringTok{"\#009E73"}\NormalTok{, }\CommentTok{\# green}
  \StringTok{"Proteobacteria"} \OtherTok{=} \StringTok{"\#F0E442"}\NormalTok{, }\CommentTok{\# yellow}
  \StringTok{"Fusobacteria"} \OtherTok{=} \StringTok{"\#CC6677"}\NormalTok{, }\CommentTok{\# dark blue}
  \StringTok{"Spirochaetae"} \OtherTok{=} \StringTok{"\#D55E00"}\NormalTok{, }\CommentTok{\# reddish{-}orange}
  \StringTok{"Cyanobacteria"} \OtherTok{=} \StringTok{"\#CC79A7"}\NormalTok{, }\CommentTok{\# purple}
  \StringTok{"Acidobacteria"} \OtherTok{=} \StringTok{"\#999999"}\NormalTok{, }\CommentTok{\# grey}
  \StringTok{"Candidate division SR1"} \OtherTok{=} \StringTok{"\#AD7700"}\NormalTok{,}\CommentTok{\# brown}
  \StringTok{"Planctomycetes"} \OtherTok{=} \StringTok{"\#332288"}\NormalTok{, }\CommentTok{\# deep indigo}
  \StringTok{"Saccharibacteria"} \OtherTok{=} \StringTok{"\#44AA99"}\NormalTok{, }\CommentTok{\# teal}
  \StringTok{"Synergistetes"} \OtherTok{=} \StringTok{"\#88CCEE"}\NormalTok{, }\CommentTok{\# light blue}
  \StringTok{"Tenericutes"} \OtherTok{=} \StringTok{"\#117733"}\NormalTok{, }\CommentTok{\# dark green}
  \StringTok{"unclassified"} \OtherTok{=} \StringTok{"\#DDDDDD"}\NormalTok{, }\CommentTok{\# very light grey}
  \StringTok{"NA"} \OtherTok{=} \StringTok{"\#DDDDDD"}\NormalTok{, }\CommentTok{\# very light grey}
  \StringTok{"Verrucomicrobia"} \OtherTok{=} \StringTok{"\#0072B2"} \CommentTok{\# reddish pink}
\NormalTok{)}
\NormalTok{mean\_abundance\_df}\SpecialCharTok{$}\NormalTok{phylum }\OtherTok{\textless{}{-}} \FunctionTok{factor}\NormalTok{(mean\_abundance\_df}\SpecialCharTok{$}\NormalTok{phylum, }\AttributeTok{levels =} \FunctionTok{names}\NormalTok{(phylum\_colors))}

\CommentTok{\# Step 12: Enhanced parallel coordinate plot}
\NormalTok{p\_parallel }\OtherTok{\textless{}{-}} \FunctionTok{ggparcoord}\NormalTok{(}
  \AttributeTok{data =}\NormalTok{ mean\_abundance\_df,}
  \AttributeTok{columns =} \DecValTok{1}\SpecialCharTok{:}\NormalTok{(}\FunctionTok{ncol}\NormalTok{(mean\_abundance\_df) }\SpecialCharTok{{-}} \DecValTok{3}\NormalTok{),  }\CommentTok{\# Exclude genus, cluster, phylum}
  \AttributeTok{groupColumn =} \StringTok{"phylum"}\NormalTok{,}
  \AttributeTok{scale =} \StringTok{"globalminmax"}\NormalTok{,  }\CommentTok{\# Use standardized values as is}
  \AttributeTok{title =} \StringTok{"Genus Abundance Trends Across Age Groups by Cluster"}\NormalTok{,}
  \AttributeTok{alphaLines =} \FloatTok{0.5}\NormalTok{,  }\CommentTok{\# Reduce overplotting}
  \AttributeTok{showPoints =} \ConstantTok{FALSE}\NormalTok{,  }\CommentTok{\# Add points for clarity}
  \AttributeTok{splineFactor =} \DecValTok{20}  \CommentTok{\# Smooth lines}
\NormalTok{) }\SpecialCharTok{+}
  \FunctionTok{facet\_wrap}\NormalTok{(}\SpecialCharTok{\textasciitilde{}}\NormalTok{ cluster, }\AttributeTok{scales =} \StringTok{"free\_y"}\NormalTok{, }\AttributeTok{ncol =} \DecValTok{2}\NormalTok{) }\SpecialCharTok{+}  \CommentTok{\# Free y{-}axis per cluster, 2 columns}
  \FunctionTok{scale\_color\_manual}\NormalTok{(}\AttributeTok{values =}\NormalTok{ phylum\_colors) }\SpecialCharTok{+}
\NormalTok{  egg}\SpecialCharTok{::}\FunctionTok{theme\_article}\NormalTok{() }\SpecialCharTok{+}  \CommentTok{\# White background, minimal grid}
  \FunctionTok{theme}\NormalTok{(}
    \AttributeTok{plot.title =} \FunctionTok{element\_text}\NormalTok{(}\AttributeTok{size =} \DecValTok{6}\NormalTok{, }\AttributeTok{face =} \StringTok{"bold"}\NormalTok{, }\AttributeTok{hjust =} \FloatTok{0.5}\NormalTok{),}
    \AttributeTok{axis.text.x =} \FunctionTok{element\_text}\NormalTok{(}\AttributeTok{angle =} \DecValTok{45}\NormalTok{, }\AttributeTok{hjust =} \DecValTok{1}\NormalTok{, }\AttributeTok{size =} \DecValTok{6}\NormalTok{),}
    \AttributeTok{axis.text.y =} \FunctionTok{element\_text}\NormalTok{(}\AttributeTok{size =} \DecValTok{6}\NormalTok{),}
    \AttributeTok{axis.title =} \FunctionTok{element\_text}\NormalTok{(}\AttributeTok{size =} \DecValTok{6}\NormalTok{),}
    \AttributeTok{strip.text =} \FunctionTok{element\_text}\NormalTok{(}\AttributeTok{size =} \DecValTok{6}\NormalTok{, }\AttributeTok{face =} \StringTok{"bold"}\NormalTok{),}
    \AttributeTok{legend.position =} \StringTok{"right"}\NormalTok{,}
    \AttributeTok{legend.title =} \FunctionTok{element\_text}\NormalTok{(}\AttributeTok{size =} \DecValTok{6}\NormalTok{, }\AttributeTok{face =} \StringTok{"bold"}\NormalTok{),}
    \AttributeTok{legend.text =} \FunctionTok{element\_text}\NormalTok{(}\AttributeTok{size =} \DecValTok{6}\NormalTok{)}
\NormalTok{  ) }\SpecialCharTok{+}
  \FunctionTok{labs}\NormalTok{(}\AttributeTok{x =} \StringTok{"Age Group"}\NormalTok{, }\AttributeTok{y =} \StringTok{"Standardized Relative Abundance"}\NormalTok{, }\AttributeTok{color =} \StringTok{"Phylum"}\NormalTok{)}

\FunctionTok{print}\NormalTok{(p\_parallel)}
\end{Highlighting}
\end{Shaded}

\pandocbounded{\includegraphics[keepaspectratio]{5_age_analyses_final_improved_files/figure-latex/unnamed-chunk-4-1.pdf}}

\begin{Shaded}
\begin{Highlighting}[]
\FunctionTok{cat}\NormalTok{(}\StringTok{"Generated enhanced parallel coordinate plot.}\SpecialCharTok{\textbackslash{}n}\StringTok{"}\NormalTok{)}
\end{Highlighting}
\end{Shaded}

\begin{verbatim}
## Generated enhanced parallel coordinate plot.
\end{verbatim}

\begin{Shaded}
\begin{Highlighting}[]
\CommentTok{\# Step 13: PCA for dimensionality reduction}
\NormalTok{pca\_result }\OtherTok{\textless{}{-}} \FunctionTok{prcomp}\NormalTok{(mean\_abundance\_scaled, }\AttributeTok{scale. =} \ConstantTok{FALSE}\NormalTok{)  }\CommentTok{\# Already scaled}
\NormalTok{pca\_df }\OtherTok{\textless{}{-}} \FunctionTok{as.data.frame}\NormalTok{(pca\_result}\SpecialCharTok{$}\NormalTok{x[, }\DecValTok{1}\SpecialCharTok{:}\DecValTok{2}\NormalTok{])  }\CommentTok{\# First two PCs}
\NormalTok{pca\_df}\SpecialCharTok{$}\NormalTok{genus }\OtherTok{\textless{}{-}} \FunctionTok{rownames}\NormalTok{(mean\_abundance\_scaled)}
\NormalTok{pca\_df}\SpecialCharTok{$}\NormalTok{cluster }\OtherTok{\textless{}{-}} \FunctionTok{factor}\NormalTok{(clusters)}
\NormalTok{pca\_df}\SpecialCharTok{$}\NormalTok{phylum }\OtherTok{\textless{}{-}}\NormalTok{ mean\_abundance\_df}\SpecialCharTok{$}\NormalTok{phylum}

\CommentTok{\# Calculate variance explained}
\NormalTok{var\_explained }\OtherTok{\textless{}{-}}\NormalTok{ pca\_result}\SpecialCharTok{$}\NormalTok{sdev}\SpecialCharTok{\^{}}\DecValTok{2} \SpecialCharTok{/} \FunctionTok{sum}\NormalTok{(pca\_result}\SpecialCharTok{$}\NormalTok{sdev}\SpecialCharTok{\^{}}\DecValTok{2}\NormalTok{) }\SpecialCharTok{*} \DecValTok{100}
\NormalTok{pc1\_label }\OtherTok{\textless{}{-}} \FunctionTok{sprintf}\NormalTok{(}\StringTok{"PC1 (\%.1f\%\%)"}\NormalTok{, var\_explained[}\DecValTok{1}\NormalTok{])}
\NormalTok{pc2\_label }\OtherTok{\textless{}{-}} \FunctionTok{sprintf}\NormalTok{(}\StringTok{"PC2 (\%.1f\%\%)"}\NormalTok{, var\_explained[}\DecValTok{2}\NormalTok{])}

\CommentTok{\# PCA plot colored by cluster}
\NormalTok{p\_pca\_cluster }\OtherTok{\textless{}{-}} \FunctionTok{ggplot}\NormalTok{(pca\_df, }\FunctionTok{aes}\NormalTok{(}\AttributeTok{x =}\NormalTok{ PC1, }\AttributeTok{y =}\NormalTok{ PC2, }\AttributeTok{color =}\NormalTok{ cluster)) }\SpecialCharTok{+}
  \FunctionTok{geom\_point}\NormalTok{(}\AttributeTok{size =} \DecValTok{1}\NormalTok{, }\AttributeTok{alpha =} \FloatTok{0.8}\NormalTok{) }\SpecialCharTok{+}
  \FunctionTok{scale\_color\_manual}\NormalTok{(}\AttributeTok{values =}\NormalTok{ cluster\_colors) }\SpecialCharTok{+}
\NormalTok{  egg}\SpecialCharTok{::}\FunctionTok{theme\_article}\NormalTok{() }\SpecialCharTok{+}  \CommentTok{\# White background, minimal grid}
  \FunctionTok{labs}\NormalTok{(}
    \AttributeTok{title =} \StringTok{"PCA of Genera by Cluster Assignment"}\NormalTok{,}
    \AttributeTok{x =}\NormalTok{ pc1\_label,}
    \AttributeTok{y =}\NormalTok{ pc2\_label,}
    \AttributeTok{color =} \StringTok{"Cluster"}
\NormalTok{  ) }\SpecialCharTok{+}
  \FunctionTok{theme}\NormalTok{(}
    \AttributeTok{plot.title =} \FunctionTok{element\_text}\NormalTok{(}\AttributeTok{size =} \DecValTok{6}\NormalTok{, }\AttributeTok{face =} \StringTok{"bold"}\NormalTok{, }\AttributeTok{hjust =} \FloatTok{0.5}\NormalTok{),}
    \AttributeTok{axis.title =} \FunctionTok{element\_text}\NormalTok{(}\AttributeTok{size =} \DecValTok{6}\NormalTok{),}
    \AttributeTok{axis.text =} \FunctionTok{element\_text}\NormalTok{(}\AttributeTok{size =} \DecValTok{6}\NormalTok{),}
    \AttributeTok{legend.title =} \FunctionTok{element\_text}\NormalTok{(}\AttributeTok{size =} \DecValTok{6}\NormalTok{, }\AttributeTok{face =} \StringTok{"bold"}\NormalTok{),}
    \AttributeTok{legend.text =} \FunctionTok{element\_text}\NormalTok{(}\AttributeTok{size =} \DecValTok{6}\NormalTok{)}
\NormalTok{  )}

\FunctionTok{print}\NormalTok{(p\_pca\_cluster)}
\end{Highlighting}
\end{Shaded}

\pandocbounded{\includegraphics[keepaspectratio]{5_age_analyses_final_improved_files/figure-latex/unnamed-chunk-4-2.pdf}}

\begin{Shaded}
\begin{Highlighting}[]
\CommentTok{\# PCA plot colored by phylum}
\NormalTok{p\_pca\_phylum }\OtherTok{\textless{}{-}} \FunctionTok{ggplot}\NormalTok{(pca\_df, }\FunctionTok{aes}\NormalTok{(}\AttributeTok{x =}\NormalTok{ PC1, }\AttributeTok{y =}\NormalTok{ PC2, }\AttributeTok{color =}\NormalTok{ phylum)) }\SpecialCharTok{+}
  \FunctionTok{geom\_point}\NormalTok{(}\AttributeTok{size =} \DecValTok{1}\NormalTok{, }\AttributeTok{alpha =} \FloatTok{0.8}\NormalTok{) }\SpecialCharTok{+}
  \FunctionTok{scale\_color\_manual}\NormalTok{(}\AttributeTok{values =}\NormalTok{ phylum\_colors) }\SpecialCharTok{+}
\NormalTok{  egg}\SpecialCharTok{::}\FunctionTok{theme\_article}\NormalTok{() }\SpecialCharTok{+}  \CommentTok{\# White background, minimal grid}
  \FunctionTok{labs}\NormalTok{(}
    \AttributeTok{title =} \StringTok{"PCA of Genera by Phylum"}\NormalTok{,}
    \AttributeTok{x =}\NormalTok{ pc1\_label,}
    \AttributeTok{y =}\NormalTok{ pc2\_label,}
    \AttributeTok{color =} \StringTok{"Phylum"}
\NormalTok{  ) }\SpecialCharTok{+}
  \FunctionTok{theme}\NormalTok{(}
    \AttributeTok{plot.title =} \FunctionTok{element\_text}\NormalTok{(}\AttributeTok{size =} \DecValTok{6}\NormalTok{, }\AttributeTok{face =} \StringTok{"bold"}\NormalTok{, }\AttributeTok{hjust =} \FloatTok{0.5}\NormalTok{),}
    \AttributeTok{axis.title =} \FunctionTok{element\_text}\NormalTok{(}\AttributeTok{size =} \DecValTok{6}\NormalTok{),}
    \AttributeTok{axis.text =} \FunctionTok{element\_text}\NormalTok{(}\AttributeTok{size =} \DecValTok{6}\NormalTok{),}
    \AttributeTok{legend.title =} \FunctionTok{element\_text}\NormalTok{(}\AttributeTok{size =} \DecValTok{6}\NormalTok{, }\AttributeTok{face =} \StringTok{"bold"}\NormalTok{),}
    \AttributeTok{legend.text =} \FunctionTok{element\_text}\NormalTok{(}\AttributeTok{size =} \DecValTok{6}\NormalTok{)}
\NormalTok{  )}

\FunctionTok{print}\NormalTok{(p\_pca\_phylum)}
\end{Highlighting}
\end{Shaded}

\pandocbounded{\includegraphics[keepaspectratio]{5_age_analyses_final_improved_files/figure-latex/unnamed-chunk-4-3.pdf}}

\begin{Shaded}
\begin{Highlighting}[]
\CommentTok{\# Step 14: t{-}SNE for non{-}linear dimensionality reduction}
\FunctionTok{set.seed}\NormalTok{(}\DecValTok{42}\NormalTok{)  }\CommentTok{\# For reproducibility}
\NormalTok{tsne\_result }\OtherTok{\textless{}{-}} \FunctionTok{Rtsne}\NormalTok{(mean\_abundance\_scaled, }\AttributeTok{dims =} \DecValTok{2}\NormalTok{, }
                     \AttributeTok{perplexity =} \FunctionTok{min}\NormalTok{(}\DecValTok{30}\NormalTok{, (}\FunctionTok{nrow}\NormalTok{(mean\_abundance\_scaled) }\SpecialCharTok{{-}} \DecValTok{1}\NormalTok{) }\SpecialCharTok{/} \DecValTok{3}\NormalTok{), }
                     \AttributeTok{verbose =} \ConstantTok{TRUE}\NormalTok{)}
\end{Highlighting}
\end{Shaded}

\begin{verbatim}
## Performing PCA
## Read the 77 x 11 data matrix successfully!
## Using no_dims = 2, perplexity = 25.333333, and theta = 0.500000
## Computing input similarities...
## Building tree...
## Done in 0.00 seconds (sparsity = 0.987013)!
## Learning embedding...
## Iteration 50: error is 49.053853 (50 iterations in 0.00 seconds)
## Iteration 100: error is 52.428653 (50 iterations in 0.00 seconds)
## Iteration 150: error is 50.671014 (50 iterations in 0.00 seconds)
## Iteration 200: error is 50.839913 (50 iterations in 0.00 seconds)
## Iteration 250: error is 50.254495 (50 iterations in 0.00 seconds)
## Iteration 300: error is 1.497633 (50 iterations in 0.00 seconds)
## Iteration 350: error is 1.079041 (50 iterations in 0.00 seconds)
## Iteration 400: error is 0.624675 (50 iterations in 0.00 seconds)
## Iteration 450: error is 0.411003 (50 iterations in 0.00 seconds)
## Iteration 500: error is 0.158988 (50 iterations in 0.00 seconds)
## Iteration 550: error is 0.130469 (50 iterations in 0.00 seconds)
## Iteration 600: error is 0.129475 (50 iterations in 0.00 seconds)
## Iteration 650: error is 0.130722 (50 iterations in 0.00 seconds)
## Iteration 700: error is 0.130171 (50 iterations in 0.00 seconds)
## Iteration 750: error is 0.129526 (50 iterations in 0.00 seconds)
## Iteration 800: error is 0.132603 (50 iterations in 0.00 seconds)
## Iteration 850: error is 0.130714 (50 iterations in 0.00 seconds)
## Iteration 900: error is 0.130631 (50 iterations in 0.00 seconds)
## Iteration 950: error is 0.129519 (50 iterations in 0.00 seconds)
## Iteration 1000: error is 0.130020 (50 iterations in 0.00 seconds)
## Fitting performed in 0.06 seconds.
\end{verbatim}

\begin{Shaded}
\begin{Highlighting}[]
\NormalTok{tsne\_df }\OtherTok{\textless{}{-}} \FunctionTok{as.data.frame}\NormalTok{(tsne\_result}\SpecialCharTok{$}\NormalTok{Y)}
\FunctionTok{colnames}\NormalTok{(tsne\_df) }\OtherTok{\textless{}{-}} \FunctionTok{c}\NormalTok{(}\StringTok{"Dim1"}\NormalTok{, }\StringTok{"Dim2"}\NormalTok{)}
\NormalTok{tsne\_df}\SpecialCharTok{$}\NormalTok{genus }\OtherTok{\textless{}{-}} \FunctionTok{rownames}\NormalTok{(mean\_abundance\_scaled)}
\NormalTok{tsne\_df}\SpecialCharTok{$}\NormalTok{cluster }\OtherTok{\textless{}{-}} \FunctionTok{factor}\NormalTok{(clusters)}
\NormalTok{tsne\_df}\SpecialCharTok{$}\NormalTok{phylum }\OtherTok{\textless{}{-}}\NormalTok{ mean\_abundance\_df}\SpecialCharTok{$}\NormalTok{phylum}

\CommentTok{\# t{-}SNE plot colored by cluster}
\NormalTok{p\_tsne\_cluster }\OtherTok{\textless{}{-}} \FunctionTok{ggplot}\NormalTok{(tsne\_df, }\FunctionTok{aes}\NormalTok{(}\AttributeTok{x =}\NormalTok{ Dim1, }\AttributeTok{y =}\NormalTok{ Dim2, }\AttributeTok{color =}\NormalTok{ cluster)) }\SpecialCharTok{+}
  \FunctionTok{geom\_point}\NormalTok{(}\AttributeTok{size =} \DecValTok{1}\NormalTok{, }\AttributeTok{alpha =} \FloatTok{0.8}\NormalTok{) }\SpecialCharTok{+}
  \FunctionTok{scale\_color\_manual}\NormalTok{(}\AttributeTok{values =}\NormalTok{ cluster\_colors) }\SpecialCharTok{+}
\NormalTok{  egg}\SpecialCharTok{::}\FunctionTok{theme\_article}\NormalTok{() }\SpecialCharTok{+}  \CommentTok{\# White background, minimal grid}
  \FunctionTok{labs}\NormalTok{(}
    \AttributeTok{title =} \StringTok{"t{-}SNE of Genera by Cluster Assignment"}\NormalTok{,}
    \AttributeTok{x =} \StringTok{"t{-}SNE Dimension 1"}\NormalTok{,}
    \AttributeTok{y =} \StringTok{"t{-}SNE Dimension 2"}\NormalTok{,}
    \AttributeTok{color =} \StringTok{"Cluster"}
\NormalTok{  ) }\SpecialCharTok{+}
  \FunctionTok{theme}\NormalTok{(}
    \AttributeTok{plot.title =} \FunctionTok{element\_text}\NormalTok{(}\AttributeTok{size =} \DecValTok{6}\NormalTok{, }\AttributeTok{face =} \StringTok{"bold"}\NormalTok{, }\AttributeTok{hjust =} \FloatTok{0.5}\NormalTok{),}
    \AttributeTok{axis.title =} \FunctionTok{element\_text}\NormalTok{(}\AttributeTok{size =} \DecValTok{6}\NormalTok{),}
    \AttributeTok{axis.text =} \FunctionTok{element\_text}\NormalTok{(}\AttributeTok{size =} \DecValTok{6}\NormalTok{),}
    \AttributeTok{legend.title =} \FunctionTok{element\_text}\NormalTok{(}\AttributeTok{size =} \DecValTok{6}\NormalTok{, }\AttributeTok{face =} \StringTok{"bold"}\NormalTok{),}
    \AttributeTok{legend.text =} \FunctionTok{element\_text}\NormalTok{(}\AttributeTok{size =} \DecValTok{6}\NormalTok{)}
\NormalTok{  )}

\FunctionTok{print}\NormalTok{(p\_tsne\_cluster)}
\end{Highlighting}
\end{Shaded}

\pandocbounded{\includegraphics[keepaspectratio]{5_age_analyses_final_improved_files/figure-latex/unnamed-chunk-4-4.pdf}}

\begin{Shaded}
\begin{Highlighting}[]
\CommentTok{\# t{-}SNE plot colored by phylum}
\NormalTok{p\_tsne\_phylum }\OtherTok{\textless{}{-}} \FunctionTok{ggplot}\NormalTok{(tsne\_df, }\FunctionTok{aes}\NormalTok{(}\AttributeTok{x =}\NormalTok{ Dim1, }\AttributeTok{y =}\NormalTok{ Dim2, }\AttributeTok{color =}\NormalTok{ phylum)) }\SpecialCharTok{+}
  \FunctionTok{geom\_point}\NormalTok{(}\AttributeTok{size =} \DecValTok{1}\NormalTok{, }\AttributeTok{alpha =} \FloatTok{0.8}\NormalTok{) }\SpecialCharTok{+}
  \FunctionTok{scale\_color\_manual}\NormalTok{(}\AttributeTok{values =}\NormalTok{ phylum\_colors) }\SpecialCharTok{+}
\NormalTok{  egg}\SpecialCharTok{::}\FunctionTok{theme\_article}\NormalTok{() }\SpecialCharTok{+}  \CommentTok{\# White background, minimal grid}
  \FunctionTok{labs}\NormalTok{(}
    \AttributeTok{title =} \StringTok{"t{-}SNE of Genera by Phylum"}\NormalTok{,}
    \AttributeTok{x =} \StringTok{"t{-}SNE Dimension 1"}\NormalTok{,}
    \AttributeTok{y =} \StringTok{"t{-}SNE Dimension 2"}\NormalTok{,}
    \AttributeTok{color =} \StringTok{"Phylum"}
\NormalTok{  ) }\SpecialCharTok{+}
  \FunctionTok{theme}\NormalTok{(}
    \AttributeTok{plot.title =} \FunctionTok{element\_text}\NormalTok{(}\AttributeTok{size =} \DecValTok{6}\NormalTok{, }\AttributeTok{face =} \StringTok{"bold"}\NormalTok{, }\AttributeTok{hjust =} \FloatTok{0.5}\NormalTok{),}
    \AttributeTok{axis.title =} \FunctionTok{element\_text}\NormalTok{(}\AttributeTok{size =} \DecValTok{6}\NormalTok{),}
    \AttributeTok{axis.text =} \FunctionTok{element\_text}\NormalTok{(}\AttributeTok{size =} \DecValTok{6}\NormalTok{),}
    \AttributeTok{legend.title =} \FunctionTok{element\_text}\NormalTok{(}\AttributeTok{size =} \DecValTok{6}\NormalTok{, }\AttributeTok{face =} \StringTok{"bold"}\NormalTok{),}
    \AttributeTok{legend.text =} \FunctionTok{element\_text}\NormalTok{(}\AttributeTok{size =} \DecValTok{6}\NormalTok{)}
\NormalTok{  )}

\FunctionTok{print}\NormalTok{(p\_tsne\_phylum)}
\end{Highlighting}
\end{Shaded}

\pandocbounded{\includegraphics[keepaspectratio]{5_age_analyses_final_improved_files/figure-latex/unnamed-chunk-4-5.pdf}}

\begin{Shaded}
\begin{Highlighting}[]
\CommentTok{\# Step 15: Create meaningful labels for heatmap rows}
\CommentTok{\# Extract taxonomy table from phyloseq object and add OTU as a column}
\NormalTok{tax\_df }\OtherTok{\textless{}{-}} \FunctionTok{as.data.frame}\NormalTok{(}\FunctionTok{tax\_table}\NormalTok{(ubiome\_filtered))}
\NormalTok{tax\_df}\SpecialCharTok{$}\NormalTok{OTU }\OtherTok{\textless{}{-}} \FunctionTok{rownames}\NormalTok{(tax\_df)  }\CommentTok{\# Convert row names (OTU IDs) to a column}

\CommentTok{\# Join taxonomy data with mean\_abundance\_df and create labels}
\NormalTok{mean\_abundance\_df }\OtherTok{\textless{}{-}}\NormalTok{ mean\_abundance\_df }\SpecialCharTok{\%\textgreater{}\%}
  \FunctionTok{left\_join}\NormalTok{(tax\_df, }\AttributeTok{by =} \FunctionTok{c}\NormalTok{(}\StringTok{"genus"} \OtherTok{=} \StringTok{"OTU"}\NormalTok{)) }\SpecialCharTok{\%\textgreater{}\%}
  \CommentTok{\# Clean the OTU ID by removing "RSV\_" and "\_relative"}
  \FunctionTok{mutate}\NormalTok{(}\AttributeTok{cleaned\_otu =} \FunctionTok{gsub}\NormalTok{(}\StringTok{"\^{}RSV\_|\^{}OTU\_"}\NormalTok{, }\StringTok{""}\NormalTok{, genus)) }\SpecialCharTok{\%\textgreater{}\%}
  \FunctionTok{mutate}\NormalTok{(}\AttributeTok{cleaned\_otu =} \FunctionTok{gsub}\NormalTok{(}\StringTok{"\_relative$"}\NormalTok{, }\StringTok{""}\NormalTok{, cleaned\_otu)) }\SpecialCharTok{\%\textgreater{}\%}
  \CommentTok{\# Create labels: use Genus if valid, otherwise "NA; \textless{}cleaned\_otu\textgreater{}"}
  \FunctionTok{mutate}\NormalTok{(}\AttributeTok{label =} \FunctionTok{if\_else}\NormalTok{(}
    \FunctionTok{is.na}\NormalTok{(Genus) }\SpecialCharTok{|} \FunctionTok{tolower}\NormalTok{(Genus) }\SpecialCharTok{==} \StringTok{"unclassified"}\NormalTok{,}
    \FunctionTok{paste0}\NormalTok{(}\StringTok{"NA; "}\NormalTok{, cleaned\_otu),}
    \FunctionTok{as.character}\NormalTok{(Genus)}
\NormalTok{  )) }\SpecialCharTok{\%\textgreater{}\%}
  \CommentTok{\# Ensure unique labels by appending OTU ID if duplicates exist}
  \FunctionTok{group\_by}\NormalTok{(label) }\SpecialCharTok{\%\textgreater{}\%}
  \FunctionTok{mutate}\NormalTok{(}\AttributeTok{label =} \FunctionTok{if\_else}\NormalTok{(}\FunctionTok{n}\NormalTok{() }\SpecialCharTok{\textgreater{}} \DecValTok{1}\NormalTok{, }\FunctionTok{paste0}\NormalTok{(label, }\StringTok{"\_"}\NormalTok{, genus), label)) }\SpecialCharTok{\%\textgreater{}\%}
  \FunctionTok{ungroup}\NormalTok{()}

\CommentTok{\# Update row names of the scaled abundance matrix with the new labels}
\FunctionTok{rownames}\NormalTok{(mean\_abundance\_scaled) }\OtherTok{\textless{}{-}}\NormalTok{ mean\_abundance\_df}\SpecialCharTok{$}\NormalTok{label  }\CommentTok{\# Use labels instead of genus}

\CommentTok{\# Step 16: Heatmap with cluster and phylum annotations}
\NormalTok{annotation\_row }\OtherTok{\textless{}{-}} \FunctionTok{data.frame}\NormalTok{(}
  \AttributeTok{Cluster =} \FunctionTok{factor}\NormalTok{(clusters),}
  \AttributeTok{Phylum =}\NormalTok{ mean\_abundance\_df}\SpecialCharTok{$}\NormalTok{phylum}
\NormalTok{)}
\FunctionTok{rownames}\NormalTok{(annotation\_row) }\OtherTok{\textless{}{-}}\NormalTok{ mean\_abundance\_df}\SpecialCharTok{$}\NormalTok{label  }\CommentTok{\# Use labels instead of genus}

\FunctionTok{pheatmap}\NormalTok{(}
\NormalTok{  mean\_abundance\_scaled,}
  \AttributeTok{cluster\_rows =}\NormalTok{ hc,}
  \AttributeTok{cluster\_cols =} \ConstantTok{FALSE}\NormalTok{,}
  \AttributeTok{annotation\_row =}\NormalTok{ annotation\_row,}
  \AttributeTok{annotation\_colors =} \FunctionTok{list}\NormalTok{(}\AttributeTok{Phylum =}\NormalTok{ phylum\_colors),}
  \AttributeTok{show\_rownames =} \ConstantTok{TRUE}\NormalTok{,  }\CommentTok{\# Display the new labels}
  \AttributeTok{color =} \FunctionTok{colorRampPalette}\NormalTok{(}\FunctionTok{c}\NormalTok{(}\StringTok{"\#053061"}\NormalTok{, }\StringTok{"\#F7F7F7"}\NormalTok{, }\StringTok{"\#67001F"}\NormalTok{))(}\DecValTok{100}\NormalTok{),}
  \AttributeTok{main =} \StringTok{"Heatmap of Standardized Genus Abundance by Age Group"}\NormalTok{,}
  \AttributeTok{fontsize =} \DecValTok{6}\NormalTok{,}
  \AttributeTok{border\_color =} \ConstantTok{NA}\NormalTok{,}
  \AttributeTok{cellwidth =} \DecValTok{10}\NormalTok{,  }\CommentTok{\# Adjust cell width for better readability}
  \AttributeTok{cellheight =} \DecValTok{10}  \CommentTok{\# Adjust cell height to fit labels}
\NormalTok{)}
\end{Highlighting}
\end{Shaded}

\pandocbounded{\includegraphics[keepaspectratio]{5_age_analyses_final_improved_files/figure-latex/unnamed-chunk-4-6.pdf}}

\begin{Shaded}
\begin{Highlighting}[]
\FunctionTok{cat}\NormalTok{(}\StringTok{"Generated heatmap with cluster and phylum annotations.}\SpecialCharTok{\textbackslash{}n}\StringTok{"}\NormalTok{)}
\end{Highlighting}
\end{Shaded}

\begin{verbatim}
## Generated heatmap with cluster and phylum annotations.
\end{verbatim}

\begin{Shaded}
\begin{Highlighting}[]
\CommentTok{\# Step 17: Cluster summary statistics}
\NormalTok{cluster\_summary }\OtherTok{\textless{}{-}}\NormalTok{ mean\_abundance\_df }\SpecialCharTok{\%\textgreater{}\%}
  \FunctionTok{group\_by}\NormalTok{(cluster) }\SpecialCharTok{\%\textgreater{}\%}
  \FunctionTok{summarise}\NormalTok{(}
    \AttributeTok{n\_genera =} \FunctionTok{n}\NormalTok{(),}
    \AttributeTok{mean\_abundance =} \FunctionTok{mean}\NormalTok{(}\FunctionTok{rowMeans}\NormalTok{(mean\_abundance\_scaled)),}
    \AttributeTok{sd\_abundance =} \FunctionTok{sd}\NormalTok{(}\FunctionTok{rowMeans}\NormalTok{(mean\_abundance\_scaled)),}
    \AttributeTok{dominant\_phyla =} \FunctionTok{paste}\NormalTok{(}\FunctionTok{names}\NormalTok{(}\FunctionTok{sort}\NormalTok{(}\FunctionTok{table}\NormalTok{(phylum), }\AttributeTok{decreasing =} \ConstantTok{TRUE}\NormalTok{))[}\DecValTok{1}\SpecialCharTok{:}\DecValTok{2}\NormalTok{], }\AttributeTok{collapse =} \StringTok{", "}\NormalTok{)}
\NormalTok{  )}

\FunctionTok{print}\NormalTok{(cluster\_summary)}
\end{Highlighting}
\end{Shaded}

\begin{verbatim}
## # A tibble: 6 x 5
##   cluster n_genera mean_abundance sd_abundance dominant_phyla            
##   <fct>      <int>          <dbl>        <dbl> <chr>                     
## 1 a             22       5.59e-17     9.08e-16 Firmicutes, Bacteroidetes 
## 2 b             14       5.59e-17     9.08e-16 Firmicutes, Actinobacteria
## 3 c              5       5.59e-17     9.08e-16 Proteobacteria, Firmicutes
## 4 d             13       5.59e-17     9.08e-16 Proteobacteria, Firmicutes
## 5 e              7       5.59e-17     9.08e-16 Actinobacteria, Firmicutes
## 6 f             16       5.59e-17     9.08e-16 Firmicutes, Proteobacteria
\end{verbatim}

\begin{Shaded}
\begin{Highlighting}[]
\CommentTok{\# Step 18: Combine PCA and t{-}SNE plots for publication}
\NormalTok{combined\_plot }\OtherTok{\textless{}{-}} \FunctionTok{plot\_grid}\NormalTok{(}
\NormalTok{  p\_pca\_cluster }\SpecialCharTok{+} \FunctionTok{theme}\NormalTok{(}\AttributeTok{legend.position =} \StringTok{"bottom"}\NormalTok{),}
\NormalTok{  p\_pca\_phylum }\SpecialCharTok{+} \FunctionTok{theme}\NormalTok{(}\AttributeTok{legend.position =} \StringTok{"bottom"}\NormalTok{),}
\NormalTok{  p\_tsne\_cluster }\SpecialCharTok{+} \FunctionTok{theme}\NormalTok{(}\AttributeTok{legend.position =} \StringTok{"bottom"}\NormalTok{),}
\NormalTok{  p\_tsne\_phylum }\SpecialCharTok{+} \FunctionTok{theme}\NormalTok{(}\AttributeTok{legend.position =} \StringTok{"bottom"}\NormalTok{),}
  \AttributeTok{labels =} \FunctionTok{c}\NormalTok{(}\StringTok{"A"}\NormalTok{, }\StringTok{"B"}\NormalTok{, }\StringTok{"C"}\NormalTok{, }\StringTok{"D"}\NormalTok{),}
  \AttributeTok{ncol =} \DecValTok{2}\NormalTok{,}
  \AttributeTok{label\_size =} \DecValTok{14}
\NormalTok{)}

\FunctionTok{print}\NormalTok{(combined\_plot)}
\end{Highlighting}
\end{Shaded}

\pandocbounded{\includegraphics[keepaspectratio]{5_age_analyses_final_improved_files/figure-latex/unnamed-chunk-4-7.pdf}}

\begin{Shaded}
\begin{Highlighting}[]
\CommentTok{\# Step 19: Save all plots and results using the PDF saving function}
\FunctionTok{save\_pdf\_figure}\NormalTok{(}\StringTok{"15.3\_parallel\_coord\_plot\_faceted\_improved.pdf"}\NormalTok{, }\AttributeTok{width\_mm =} \DecValTok{100}\NormalTok{, }\AttributeTok{height\_mm =} \DecValTok{90}\NormalTok{)}
\FunctionTok{print}\NormalTok{(p\_parallel)}
\FunctionTok{dev.off}\NormalTok{()}
\end{Highlighting}
\end{Shaded}

\begin{verbatim}
## pdf 
##   2
\end{verbatim}

\begin{Shaded}
\begin{Highlighting}[]
\FunctionTok{save\_pdf\_figure}\NormalTok{(}\StringTok{"15.4\_pca\_cluster\_improved.pdf"}\NormalTok{, }\AttributeTok{width\_mm =} \DecValTok{44}\NormalTok{, }\AttributeTok{height\_mm =} \DecValTok{30}\NormalTok{)}
\FunctionTok{print}\NormalTok{(p\_pca\_cluster)}
\FunctionTok{dev.off}\NormalTok{()}
\end{Highlighting}
\end{Shaded}

\begin{verbatim}
## pdf 
##   2
\end{verbatim}

\begin{Shaded}
\begin{Highlighting}[]
\FunctionTok{save\_pdf\_figure}\NormalTok{(}\StringTok{"15.5\_pca\_phylum\_improved.pdf"}\NormalTok{, }\AttributeTok{width\_mm =} \DecValTok{65}\NormalTok{, }\AttributeTok{height\_mm =} \DecValTok{30}\NormalTok{)}
\FunctionTok{print}\NormalTok{(p\_pca\_phylum)}
\FunctionTok{dev.off}\NormalTok{()}
\end{Highlighting}
\end{Shaded}

\begin{verbatim}
## pdf 
##   2
\end{verbatim}

\begin{Shaded}
\begin{Highlighting}[]
\FunctionTok{save\_pdf\_figure}\NormalTok{(}\StringTok{"15.6\_tsne\_cluster\_improved.pdf"}\NormalTok{, }\AttributeTok{width\_mm =} \DecValTok{44}\NormalTok{, }\AttributeTok{height\_mm =} \DecValTok{30}\NormalTok{)}
\FunctionTok{print}\NormalTok{(p\_tsne\_cluster)}
\FunctionTok{dev.off}\NormalTok{()}
\end{Highlighting}
\end{Shaded}

\begin{verbatim}
## pdf 
##   2
\end{verbatim}

\begin{Shaded}
\begin{Highlighting}[]
\FunctionTok{save\_pdf\_figure}\NormalTok{(}\StringTok{"15.7\_tsne\_phylum\_improved.pdf"}\NormalTok{, }\AttributeTok{width\_mm =} \DecValTok{65}\NormalTok{, }\AttributeTok{height\_mm =} \DecValTok{30}\NormalTok{)}
\FunctionTok{print}\NormalTok{(p\_tsne\_phylum)}
\FunctionTok{dev.off}\NormalTok{()}
\end{Highlighting}
\end{Shaded}

\begin{verbatim}
## pdf 
##   2
\end{verbatim}

\begin{Shaded}
\begin{Highlighting}[]
\FunctionTok{save\_pdf\_figure}\NormalTok{(}\StringTok{"15.8\_combined\_pca\_tsne\_improved.pdf"}\NormalTok{, }\AttributeTok{width\_mm =} \DecValTok{65}\NormalTok{, }\AttributeTok{height\_mm =} \DecValTok{100}\NormalTok{)}
\FunctionTok{print}\NormalTok{(combined\_plot)}
\FunctionTok{dev.off}\NormalTok{()}
\end{Highlighting}
\end{Shaded}

\begin{verbatim}
## pdf 
##   2
\end{verbatim}

\begin{Shaded}
\begin{Highlighting}[]
\FunctionTok{save\_pdf\_figure}\NormalTok{(}\StringTok{"15.9\_heatmap\_improved.pdf"}\NormalTok{, }\AttributeTok{width\_mm =} \DecValTok{110}\NormalTok{, }\AttributeTok{height\_mm =} \DecValTok{200}\NormalTok{)}
\FunctionTok{pheatmap}\NormalTok{(}
\NormalTok{  mean\_abundance\_scaled,}
  \AttributeTok{cluster\_rows =}\NormalTok{ hc,}
  \AttributeTok{cluster\_cols =} \ConstantTok{FALSE}\NormalTok{,}
  \AttributeTok{annotation\_row =}\NormalTok{ annotation\_row,}
  \AttributeTok{annotation\_colors =} \FunctionTok{list}\NormalTok{(}\AttributeTok{Phylum =}\NormalTok{ phylum\_colors),}
  \AttributeTok{show\_rownames =} \ConstantTok{TRUE}\NormalTok{,  }\CommentTok{\# Displays the new Genus{-}based labels}
  \AttributeTok{color =} \FunctionTok{colorRampPalette}\NormalTok{(}\FunctionTok{c}\NormalTok{(}\StringTok{"\#053061"}\NormalTok{, }\StringTok{"\#F7F7F7"}\NormalTok{, }\StringTok{"\#67001F"}\NormalTok{))(}\DecValTok{100}\NormalTok{),}
  \AttributeTok{main =} \StringTok{"Heatmap of Standardized Genus Abundance by Age Group"}\NormalTok{,}
  \AttributeTok{fontsize =} \DecValTok{6}\NormalTok{,}
  \AttributeTok{border\_color =} \ConstantTok{NA}\NormalTok{,}
  \AttributeTok{cellwidth =} \DecValTok{5}\NormalTok{,}
  \AttributeTok{cellheight =} \DecValTok{5}
\NormalTok{)}
\FunctionTok{dev.off}\NormalTok{()}
\end{Highlighting}
\end{Shaded}

\begin{verbatim}
## pdf 
##   2
\end{verbatim}

\begin{Shaded}
\begin{Highlighting}[]
\CommentTok{\# Save data frames as CSV}
\FunctionTok{write.csv}\NormalTok{(mean\_abundance\_df, }\FunctionTok{file.path}\NormalTok{(viz\_out\_path, }\StringTok{"genus\_cluster\_assignments\_improved.csv"}\NormalTok{), }\AttributeTok{row.names =} \ConstantTok{FALSE}\NormalTok{)}
\FunctionTok{write.csv}\NormalTok{(cluster\_summary, }\FunctionTok{file.path}\NormalTok{(viz\_out\_path, }\StringTok{"cluster\_summary\_improved.csv"}\NormalTok{), }\AttributeTok{row.names =} \ConstantTok{FALSE}\NormalTok{)}

\FunctionTok{cat}\NormalTok{(}\StringTok{"Saved all plots and results.}\SpecialCharTok{\textbackslash{}n}\StringTok{"}\NormalTok{)}
\end{Highlighting}
\end{Shaded}

\begin{verbatim}
## Saved all plots and results.
\end{verbatim}

\subsubsection{Complete R Code for Additional
Analysis}\label{complete-r-code-for-additional-analysis}

\begin{Shaded}
\begin{Highlighting}[]
\CommentTok{\# Set seed for reproducibility}
\FunctionTok{set.seed}\NormalTok{(}\DecValTok{42}\NormalTok{)}

\CommentTok{\# Required libraries}
\FunctionTok{library}\NormalTok{(phyloseq) }\CommentTok{\# For microbiome data handling}
\FunctionTok{library}\NormalTok{(dplyr) }\CommentTok{\# For data manipulation}
\FunctionTok{library}\NormalTok{(GGally) }\CommentTok{\# For parallel coordinate plots}
\FunctionTok{library}\NormalTok{(cluster) }\CommentTok{\# For silhouette analysis}
\FunctionTok{library}\NormalTok{(Rtsne) }\CommentTok{\# For t{-}SNE}
\FunctionTok{library}\NormalTok{(pheatmap) }\CommentTok{\# For heatmaps}
\FunctionTok{library}\NormalTok{(ggplot2) }\CommentTok{\# For plotting}
\FunctionTok{library}\NormalTok{(cowplot) }\CommentTok{\# For arranging multiple plots}
\FunctionTok{library}\NormalTok{(egg) }\CommentTok{\# For theme\_article()}

\CommentTok{\# Function to save PDF figures with standardized settings}
\NormalTok{save\_pdf\_figure }\OtherTok{\textless{}{-}} \ControlFlowTok{function}\NormalTok{(file\_name, }\AttributeTok{width\_mm =} \DecValTok{180}\NormalTok{, }\AttributeTok{height\_mm =} \DecValTok{170}\NormalTok{) \{}
  \FunctionTok{pdf}\NormalTok{(}
    \AttributeTok{file =} \FunctionTok{file.path}\NormalTok{(viz\_out\_path, file\_name),}
    \AttributeTok{width =}\NormalTok{ width\_mm }\SpecialCharTok{/} \FloatTok{25.4}\NormalTok{,  }\CommentTok{\# Convert mm to inches}
    \AttributeTok{height =}\NormalTok{ height\_mm }\SpecialCharTok{/} \FloatTok{25.4}\NormalTok{,}
    \AttributeTok{family =} \StringTok{"Helvetica"}\NormalTok{,     }\CommentTok{\# Use Helvetica font}
    \AttributeTok{colormodel =} \StringTok{"rgb"}\NormalTok{,       }\CommentTok{\# Ensure color consistency}
    \AttributeTok{pointsize =} \DecValTok{5}\NormalTok{,            }\CommentTok{\# Default base point size}
    \AttributeTok{useDingbats =} \ConstantTok{FALSE}       \CommentTok{\# Prevent font embedding issues}
\NormalTok{  )}
\NormalTok{\}}

\CommentTok{\# Step 4: Compute mean relative abundance per age group}
\NormalTok{otu\_mat }\OtherTok{\textless{}{-}} \FunctionTok{as}\NormalTok{(}\FunctionTok{otu\_table}\NormalTok{(ubiome\_filtered), }\StringTok{"matrix"}\NormalTok{)  }\CommentTok{\# Taxa as rows, samples as columns}
\NormalTok{sample\_df }\OtherTok{\textless{}{-}} \FunctionTok{as}\NormalTok{(}\FunctionTok{sample\_data}\NormalTok{(ubiome\_filtered), }\StringTok{"data.frame"}\NormalTok{)}
\NormalTok{otu\_t }\OtherTok{\textless{}{-}} \FunctionTok{t}\NormalTok{(otu\_mat)  }\CommentTok{\# Transpose so samples are rows, taxa are columns}
\NormalTok{otu\_df }\OtherTok{\textless{}{-}} \FunctionTok{as.data.frame}\NormalTok{(otu\_t)}
\NormalTok{otu\_df}\SpecialCharTok{$}\NormalTok{age\_group }\OtherTok{\textless{}{-}}\NormalTok{ sample\_df}\SpecialCharTok{$}\NormalTok{age\_group}

\NormalTok{mean\_abundance }\OtherTok{\textless{}{-}}\NormalTok{ otu\_df }\SpecialCharTok{\%\textgreater{}\%}
  \FunctionTok{group\_by}\NormalTok{(age\_group) }\SpecialCharTok{\%\textgreater{}\%}
  \FunctionTok{summarise}\NormalTok{(}\FunctionTok{across}\NormalTok{(}\FunctionTok{everything}\NormalTok{(), mean, }\AttributeTok{na.rm =} \ConstantTok{TRUE}\NormalTok{))}

\NormalTok{mean\_abundance\_mat }\OtherTok{\textless{}{-}} \FunctionTok{as.matrix}\NormalTok{(mean\_abundance[, }\SpecialCharTok{{-}}\DecValTok{1}\NormalTok{])  }\CommentTok{\# Remove age\_group column}
\FunctionTok{rownames}\NormalTok{(mean\_abundance\_mat) }\OtherTok{\textless{}{-}}\NormalTok{ mean\_abundance}\SpecialCharTok{$}\NormalTok{age\_group}
\FunctionTok{cat}\NormalTok{(}\StringTok{"Computed mean relative abundance per age group. Dimensions:"}\NormalTok{, }\FunctionTok{dim}\NormalTok{(mean\_abundance\_mat), }\StringTok{"}\SpecialCharTok{\textbackslash{}n}\StringTok{"}\NormalTok{)}
\end{Highlighting}
\end{Shaded}

\begin{verbatim}
## Computed mean relative abundance per age group. Dimensions: 11 93
\end{verbatim}

\begin{Shaded}
\begin{Highlighting}[]
\CommentTok{\# Step 5: Identify differentially abundant genera across age groups using Kruskal{-}Wallis test}
\NormalTok{p\_values }\OtherTok{\textless{}{-}} \FunctionTok{apply}\NormalTok{(otu\_mat, }\DecValTok{1}\NormalTok{, }\ControlFlowTok{function}\NormalTok{(x) }\FunctionTok{kruskal.test}\NormalTok{(x }\SpecialCharTok{\textasciitilde{}}\NormalTok{ sample\_df}\SpecialCharTok{$}\NormalTok{age\_group)}\SpecialCharTok{$}\NormalTok{p.value)}
\NormalTok{adj\_p\_values }\OtherTok{\textless{}{-}} \FunctionTok{p.adjust}\NormalTok{(p\_values, }\AttributeTok{method =} \StringTok{"BH"}\NormalTok{)  }\CommentTok{\# Benjamini{-}Hochberg correction}
\NormalTok{significant\_genera }\OtherTok{\textless{}{-}} \FunctionTok{names}\NormalTok{(}\FunctionTok{which}\NormalTok{(adj\_p\_values }\SpecialCharTok{\textless{}} \FloatTok{0.05}\NormalTok{))}
\FunctionTok{cat}\NormalTok{(}\StringTok{"Identified"}\NormalTok{, }\FunctionTok{length}\NormalTok{(significant\_genera), }\StringTok{"significant genera (adjusted p \textless{} 0.05).}\SpecialCharTok{\textbackslash{}n}\StringTok{"}\NormalTok{)}
\end{Highlighting}
\end{Shaded}

\begin{verbatim}
## Identified 77 significant genera (adjusted p < 0.05).
\end{verbatim}

\begin{Shaded}
\begin{Highlighting}[]
\CommentTok{\# Step 6: Subset mean abundance matrix to significant genera}
\NormalTok{mean\_abundance\_sig }\OtherTok{\textless{}{-}}\NormalTok{ mean\_abundance\_mat[, significant\_genera, drop }\OtherTok{=} \ConstantTok{FALSE}\NormalTok{]}

\CommentTok{\# Step 7: Standardize the mean abundance profiles (mean = 0, SD = 1) for each genus}
\NormalTok{mean\_abundance\_sig\_t }\OtherTok{\textless{}{-}} \FunctionTok{t}\NormalTok{(mean\_abundance\_sig)  }\CommentTok{\# Genera as rows, age groups as columns}
\NormalTok{mean\_abundance\_scaled }\OtherTok{\textless{}{-}} \FunctionTok{t}\NormalTok{(}\FunctionTok{scale}\NormalTok{(}\FunctionTok{t}\NormalTok{(mean\_abundance\_sig\_t)))  }\CommentTok{\# Scale each genus across age groups}
\FunctionTok{cat}\NormalTok{(}\StringTok{"Standardized mean abundance profiles.}\SpecialCharTok{\textbackslash{}n}\StringTok{"}\NormalTok{)}
\end{Highlighting}
\end{Shaded}

\begin{verbatim}
## Standardized mean abundance profiles.
\end{verbatim}

\begin{Shaded}
\begin{Highlighting}[]
\CommentTok{\# Step 8: Perform hierarchical clustering using Ward\textquotesingle{}s linkage}
\NormalTok{dist\_mat }\OtherTok{\textless{}{-}} \FunctionTok{dist}\NormalTok{(mean\_abundance\_scaled, }\AttributeTok{method =} \StringTok{"euclidean"}\NormalTok{)}
\NormalTok{hc }\OtherTok{\textless{}{-}} \FunctionTok{hclust}\NormalTok{(dist\_mat, }\AttributeTok{method =} \StringTok{"ward.D2"}\NormalTok{)}
\FunctionTok{cat}\NormalTok{(}\StringTok{"Performed hierarchical clustering with Ward\textquotesingle{}s method.}\SpecialCharTok{\textbackslash{}n}\StringTok{"}\NormalTok{)}
\end{Highlighting}
\end{Shaded}

\begin{verbatim}
## Performed hierarchical clustering with Ward's method.
\end{verbatim}

\begin{Shaded}
\begin{Highlighting}[]
\CommentTok{\# Step 9: Determine optimal number of clusters using simple silhouette method}
\CommentTok{\# Use a simple approach that naturally finds 3{-}6 clusters}
\NormalTok{sil\_scores }\OtherTok{\textless{}{-}} \FunctionTok{numeric}\NormalTok{()}

\ControlFlowTok{for}\NormalTok{ (k }\ControlFlowTok{in} \DecValTok{2}\SpecialCharTok{:}\FunctionTok{min}\NormalTok{(}\DecValTok{8}\NormalTok{, }\FunctionTok{nrow}\NormalTok{(mean\_abundance\_scaled))) \{  }\CommentTok{\# Test up to 8 clusters, but not more than data points}
\NormalTok{  clusters\_temp }\OtherTok{\textless{}{-}} \FunctionTok{cutree}\NormalTok{(hc, }\AttributeTok{k =}\NormalTok{ k)}
\NormalTok{  sil }\OtherTok{\textless{}{-}} \FunctionTok{silhouette}\NormalTok{(clusters\_temp, dist\_mat)}
\NormalTok{  sil\_scores[k }\SpecialCharTok{{-}} \DecValTok{1}\NormalTok{] }\OtherTok{\textless{}{-}} \FunctionTok{mean}\NormalTok{(sil)  }\CommentTok{\# Use mean function directly}
\NormalTok{\}}

\CommentTok{\# Find the optimal k using simple silhouette scores}
\NormalTok{optimal\_k }\OtherTok{\textless{}{-}} \FunctionTok{which.max}\NormalTok{(sil\_scores) }\SpecialCharTok{+} \DecValTok{1}

\CommentTok{\# Apply simple constraints (prefer 3{-}6 clusters)}
\NormalTok{optimal\_k }\OtherTok{\textless{}{-}} \FunctionTok{min}\NormalTok{(optimal\_k, }\DecValTok{6}\NormalTok{)}
\NormalTok{optimal\_k }\OtherTok{\textless{}{-}} \FunctionTok{max}\NormalTok{(optimal\_k, }\DecValTok{3}\NormalTok{)}

\FunctionTok{cat}\NormalTok{(}\StringTok{"Simple silhouette analysis results:}\SpecialCharTok{\textbackslash{}n}\StringTok{"}\NormalTok{)}
\end{Highlighting}
\end{Shaded}

\begin{verbatim}
## Simple silhouette analysis results:
\end{verbatim}

\begin{Shaded}
\begin{Highlighting}[]
\FunctionTok{cat}\NormalTok{(}\StringTok{"  Silhouette scores (k=2{-}8):"}\NormalTok{, }\FunctionTok{round}\NormalTok{(sil\_scores, }\DecValTok{3}\NormalTok{), }\StringTok{"}\SpecialCharTok{\textbackslash{}n}\StringTok{"}\NormalTok{)}
\end{Highlighting}
\end{Shaded}

\begin{verbatim}
##   Silhouette scores (k=2-8): 1.1 1.319 1.763 2.15 2.41 2.491 2.553
\end{verbatim}

\begin{Shaded}
\begin{Highlighting}[]
\FunctionTok{cat}\NormalTok{(}\StringTok{"  Optimal k (raw):"}\NormalTok{, }\FunctionTok{which.max}\NormalTok{(sil\_scores) }\SpecialCharTok{+} \DecValTok{1}\NormalTok{, }\StringTok{"clusters}\SpecialCharTok{\textbackslash{}n}\StringTok{"}\NormalTok{)}
\end{Highlighting}
\end{Shaded}

\begin{verbatim}
##   Optimal k (raw): 8 clusters
\end{verbatim}

\begin{Shaded}
\begin{Highlighting}[]
\FunctionTok{cat}\NormalTok{(}\StringTok{"  Final optimal k (constrained 3{-}6):"}\NormalTok{, optimal\_k, }\StringTok{"clusters}\SpecialCharTok{\textbackslash{}n}\StringTok{"}\NormalTok{)}
\end{Highlighting}
\end{Shaded}

\begin{verbatim}
##   Final optimal k (constrained 3-6): 6 clusters
\end{verbatim}

\begin{Shaded}
\begin{Highlighting}[]
\CommentTok{\# Cut the tree into the optimal number of clusters}
\NormalTok{clusters\_numeric }\OtherTok{\textless{}{-}} \FunctionTok{cutree}\NormalTok{(hc, }\AttributeTok{k =}\NormalTok{ optimal\_k)}
\FunctionTok{cat}\NormalTok{(}\StringTok{"Using"}\NormalTok{, optimal\_k, }\StringTok{"clusters determined by simple silhouette analysis.}\SpecialCharTok{\textbackslash{}n}\StringTok{"}\NormalTok{)}
\end{Highlighting}
\end{Shaded}

\begin{verbatim}
## Using 6 clusters determined by simple silhouette analysis.
\end{verbatim}

\begin{Shaded}
\begin{Highlighting}[]
\CommentTok{\# Step 10: Convert numeric cluster assignments to character cluster names}
\CommentTok{\# Create dynamic cluster mapping based on optimal number of clusters}
\NormalTok{cluster\_letters }\OtherTok{\textless{}{-}}\NormalTok{ letters[}\DecValTok{1}\SpecialCharTok{:}\FunctionTok{min}\NormalTok{(optimal\_k, }\DecValTok{26}\NormalTok{)]  }\CommentTok{\# Use letters a{-}z (max 26 clusters)}
\NormalTok{cluster\_mapping }\OtherTok{\textless{}{-}} \FunctionTok{setNames}\NormalTok{(cluster\_letters, }\FunctionTok{as.character}\NormalTok{(}\DecValTok{1}\SpecialCharTok{:}\NormalTok{optimal\_k))}

\CommentTok{\# Convert numeric clusters to character clusters}
\NormalTok{clusters }\OtherTok{\textless{}{-}} \FunctionTok{factor}\NormalTok{(cluster\_mapping[}\FunctionTok{as.character}\NormalTok{(clusters\_numeric)], }
                  \AttributeTok{levels =}\NormalTok{ cluster\_letters)}

\CommentTok{\# Define cluster colors (dynamic based on number of clusters)}
\CommentTok{\# Use a predefined color palette that can handle up to 26 clusters}
\NormalTok{cluster\_color\_palette }\OtherTok{\textless{}{-}} \FunctionTok{c}\NormalTok{(}
  \StringTok{"\#FF80DA"}\NormalTok{, }\StringTok{"\#FF9289"}\NormalTok{, }\StringTok{"\#F3AC00"}\NormalTok{, }\StringTok{"\#89CC00"}\NormalTok{, }\StringTok{"\#00DCA9"}\NormalTok{, }\StringTok{"\#00D8EF"}\NormalTok{, }
  \StringTok{"\#9CB1FF"}\NormalTok{, }\StringTok{"\#D69EFF"}\NormalTok{, }\StringTok{"\#FF6B6B"}\NormalTok{, }\StringTok{"\#4ECDC4"}\NormalTok{, }\StringTok{"\#45B7D1"}\NormalTok{, }\StringTok{"\#96CEB4"}\NormalTok{,}
  \StringTok{"\#FFEAA7"}\NormalTok{, }\StringTok{"\#DDA0DD"}\NormalTok{, }\StringTok{"\#98D8C8"}\NormalTok{, }\StringTok{"\#F7DC6F"}\NormalTok{, }\StringTok{"\#BB8FCE"}\NormalTok{, }\StringTok{"\#85C1E9"}\NormalTok{,}
  \StringTok{"\#F8C471"}\NormalTok{, }\StringTok{"\#82E0AA"}\NormalTok{, }\StringTok{"\#F1948A"}\NormalTok{, }\StringTok{"\#85C1E9"}\NormalTok{, }\StringTok{"\#D7BDE2"}\NormalTok{, }\StringTok{"\#A9DFBF"}\NormalTok{,}
  \StringTok{"\#F9E79F"}\NormalTok{, }\StringTok{"\#AED6F1"}
\NormalTok{)}
\NormalTok{cluster\_colors }\OtherTok{\textless{}{-}} \FunctionTok{setNames}\NormalTok{(cluster\_color\_palette[}\DecValTok{1}\SpecialCharTok{:}\NormalTok{optimal\_k], cluster\_letters)}

\CommentTok{\# Step 11: Prepare data for parallel coordinate plot}
\NormalTok{mean\_abundance\_df }\OtherTok{\textless{}{-}} \FunctionTok{as.data.frame}\NormalTok{(mean\_abundance\_scaled)}
\NormalTok{mean\_abundance\_df}\SpecialCharTok{$}\NormalTok{genus }\OtherTok{\textless{}{-}} \FunctionTok{rownames}\NormalTok{(mean\_abundance\_df)}
\NormalTok{mean\_abundance\_df}\SpecialCharTok{$}\NormalTok{cluster }\OtherTok{\textless{}{-}}\NormalTok{ clusters}

\CommentTok{\# Extract phylum information for each genus}
\NormalTok{tax\_table\_df }\OtherTok{\textless{}{-}} \FunctionTok{as.data.frame}\NormalTok{(}\FunctionTok{tax\_table}\NormalTok{(ubiome\_filtered))}
\NormalTok{mean\_abundance\_df}\SpecialCharTok{$}\NormalTok{phylum }\OtherTok{\textless{}{-}}\NormalTok{ tax\_table\_df[mean\_abundance\_df}\SpecialCharTok{$}\NormalTok{genus, }\StringTok{"Phylum"}\NormalTok{]}

\CommentTok{\# Define phylum colors}
\NormalTok{phylum\_colors }\OtherTok{\textless{}{-}} \FunctionTok{c}\NormalTok{(}
  \StringTok{"Firmicutes"} \OtherTok{=} \StringTok{"\#E69F00"}\NormalTok{, }\CommentTok{\# orange}
  \StringTok{"Bacteroidetes"} \OtherTok{=} \StringTok{"\#56B4E9"}\NormalTok{, }\CommentTok{\# blue}
  \StringTok{"Actinobacteria"} \OtherTok{=} \StringTok{"\#009E73"}\NormalTok{, }\CommentTok{\# green}
  \StringTok{"Proteobacteria"} \OtherTok{=} \StringTok{"\#F0E442"}\NormalTok{, }\CommentTok{\# yellow}
  \StringTok{"Fusobacteria"} \OtherTok{=} \StringTok{"\#CC6677"}\NormalTok{, }\CommentTok{\# dark blue}
  \StringTok{"Spirochaetae"} \OtherTok{=} \StringTok{"\#D55E00"}\NormalTok{, }\CommentTok{\# reddish{-}orange}
  \StringTok{"Cyanobacteria"} \OtherTok{=} \StringTok{"\#CC79A7"}\NormalTok{, }\CommentTok{\# purple}
  \StringTok{"Acidobacteria"} \OtherTok{=} \StringTok{"\#999999"}\NormalTok{, }\CommentTok{\# grey}
  \StringTok{"Candidate division SR1"} \OtherTok{=} \StringTok{"\#AD7700"}\NormalTok{,}\CommentTok{\# brown}
  \StringTok{"Planctomycetes"} \OtherTok{=} \StringTok{"\#332288"}\NormalTok{, }\CommentTok{\# deep indigo}
  \StringTok{"Saccharibacteria"} \OtherTok{=} \StringTok{"\#44AA99"}\NormalTok{, }\CommentTok{\# teal}
  \StringTok{"Synergistetes"} \OtherTok{=} \StringTok{"\#88CCEE"}\NormalTok{, }\CommentTok{\# light blue}
  \StringTok{"Tenericutes"} \OtherTok{=} \StringTok{"\#117733"}\NormalTok{, }\CommentTok{\# dark green}
  \StringTok{"unclassified"} \OtherTok{=} \StringTok{"\#DDDDDD"}\NormalTok{, }\CommentTok{\# very light grey}
  \StringTok{"NA"} \OtherTok{=} \StringTok{"\#DDDDDD"}\NormalTok{, }\CommentTok{\# very light grey}
  \StringTok{"Verrucomicrobia"} \OtherTok{=} \StringTok{"\#0072B2"} \CommentTok{\# reddish pink}
\NormalTok{)}

\CommentTok{\# Step 12: Create parallel coordinate plot colored by cluster}
\NormalTok{p\_cluster }\OtherTok{\textless{}{-}} \FunctionTok{ggparcoord}\NormalTok{(}
  \AttributeTok{data =}\NormalTok{ mean\_abundance\_df,}
  \AttributeTok{columns =} \DecValTok{1}\SpecialCharTok{:}\NormalTok{(}\FunctionTok{ncol}\NormalTok{(mean\_abundance\_df) }\SpecialCharTok{{-}} \DecValTok{3}\NormalTok{),  }\CommentTok{\# Exclude genus, cluster, phylum columns}
  \AttributeTok{groupColumn =} \StringTok{"cluster"}\NormalTok{,}
  \AttributeTok{scale =} \StringTok{"globalminmax"}\NormalTok{,  }\CommentTok{\# Use standardized values as is}
  \AttributeTok{title =} \StringTok{"Parallel Coordinate Plot of Genus Abundance Trends by Age Group (Cluster)"}\NormalTok{,}
  \AttributeTok{alphaLines =} \FloatTok{0.5}  \CommentTok{\# Set transparency}
\NormalTok{) }\SpecialCharTok{+}
  \FunctionTok{scale\_color\_manual}\NormalTok{(}\AttributeTok{values =}\NormalTok{ cluster\_colors) }\SpecialCharTok{+}  \CommentTok{\# Apply custom cluster colors}
\NormalTok{  egg}\SpecialCharTok{::}\FunctionTok{theme\_article}\NormalTok{() }\SpecialCharTok{+}  \CommentTok{\# White background, minimal grid}
  \FunctionTok{theme}\NormalTok{(}
    \AttributeTok{plot.title =} \FunctionTok{element\_text}\NormalTok{(}\AttributeTok{size =} \DecValTok{6}\NormalTok{, }\AttributeTok{face =} \StringTok{"bold"}\NormalTok{, }\AttributeTok{hjust =} \FloatTok{0.5}\NormalTok{),}
    \AttributeTok{axis.text.x =} \FunctionTok{element\_text}\NormalTok{(}\AttributeTok{angle =} \DecValTok{45}\NormalTok{, }\AttributeTok{hjust =} \DecValTok{1}\NormalTok{, }\AttributeTok{size =} \DecValTok{6}\NormalTok{),}
    \AttributeTok{axis.text.y =} \FunctionTok{element\_text}\NormalTok{(}\AttributeTok{size =} \DecValTok{6}\NormalTok{),}
    \AttributeTok{axis.title =} \FunctionTok{element\_text}\NormalTok{(}\AttributeTok{size =} \DecValTok{6}\NormalTok{),}
    \AttributeTok{legend.title =} \FunctionTok{element\_text}\NormalTok{(}\AttributeTok{size =} \DecValTok{6}\NormalTok{, }\AttributeTok{face =} \StringTok{"bold"}\NormalTok{),}
    \AttributeTok{legend.text =} \FunctionTok{element\_text}\NormalTok{(}\AttributeTok{size =} \DecValTok{6}\NormalTok{)}
\NormalTok{  ) }\SpecialCharTok{+}
  \FunctionTok{labs}\NormalTok{(}\AttributeTok{x =} \StringTok{"Age Group"}\NormalTok{, }\AttributeTok{y =} \StringTok{"Standardized Relative Abundance"}\NormalTok{, }\AttributeTok{color =} \StringTok{"Cluster"}\NormalTok{)}

\FunctionTok{print}\NormalTok{(p\_cluster)}
\end{Highlighting}
\end{Shaded}

\pandocbounded{\includegraphics[keepaspectratio]{5_age_analyses_final_improved_files/figure-latex/unnamed-chunk-5-1.pdf}}

\begin{Shaded}
\begin{Highlighting}[]
\FunctionTok{cat}\NormalTok{(}\StringTok{"Generated parallel coordinate plot colored by cluster.}\SpecialCharTok{\textbackslash{}n}\StringTok{"}\NormalTok{)}
\end{Highlighting}
\end{Shaded}

\begin{verbatim}
## Generated parallel coordinate plot colored by cluster.
\end{verbatim}

\begin{Shaded}
\begin{Highlighting}[]
\CommentTok{\# Step 13: Create parallel coordinate plot colored by phylum}
\NormalTok{p\_phylum }\OtherTok{\textless{}{-}} \FunctionTok{ggparcoord}\NormalTok{(}
  \AttributeTok{data =}\NormalTok{ mean\_abundance\_df,}
  \AttributeTok{columns =} \DecValTok{1}\SpecialCharTok{:}\NormalTok{(}\FunctionTok{ncol}\NormalTok{(mean\_abundance\_df) }\SpecialCharTok{{-}} \DecValTok{3}\NormalTok{),  }\CommentTok{\# Exclude genus, cluster, phylum columns}
  \AttributeTok{groupColumn =} \StringTok{"phylum"}\NormalTok{,}
  \AttributeTok{scale =} \StringTok{"globalminmax"}\NormalTok{,  }\CommentTok{\# Use standardized values as is}
  \AttributeTok{title =} \StringTok{"Parallel Coordinate Plot of Genus Abundance Trends by Age Group (Phylum)"}\NormalTok{,}
  \AttributeTok{alphaLines =} \FloatTok{0.5}  \CommentTok{\# Set transparency}
\NormalTok{) }\SpecialCharTok{+}
  \FunctionTok{scale\_color\_manual}\NormalTok{(}\AttributeTok{values =}\NormalTok{ phylum\_colors) }\SpecialCharTok{+}  \CommentTok{\# Apply custom phylum colors}
\NormalTok{  egg}\SpecialCharTok{::}\FunctionTok{theme\_article}\NormalTok{() }\SpecialCharTok{+}  \CommentTok{\# White background, minimal grid}
  \FunctionTok{theme}\NormalTok{(}
    \AttributeTok{plot.title =} \FunctionTok{element\_text}\NormalTok{(}\AttributeTok{size =} \DecValTok{6}\NormalTok{, }\AttributeTok{face =} \StringTok{"bold"}\NormalTok{, }\AttributeTok{hjust =} \FloatTok{0.5}\NormalTok{),}
    \AttributeTok{axis.text.x =} \FunctionTok{element\_text}\NormalTok{(}\AttributeTok{angle =} \DecValTok{45}\NormalTok{, }\AttributeTok{hjust =} \DecValTok{1}\NormalTok{, }\AttributeTok{size =} \DecValTok{6}\NormalTok{),}
    \AttributeTok{axis.text.y =} \FunctionTok{element\_text}\NormalTok{(}\AttributeTok{size =} \DecValTok{6}\NormalTok{),}
    \AttributeTok{axis.title =} \FunctionTok{element\_text}\NormalTok{(}\AttributeTok{size =} \DecValTok{6}\NormalTok{),}
    \AttributeTok{legend.title =} \FunctionTok{element\_text}\NormalTok{(}\AttributeTok{size =} \DecValTok{6}\NormalTok{, }\AttributeTok{face =} \StringTok{"bold"}\NormalTok{),}
    \AttributeTok{legend.text =} \FunctionTok{element\_text}\NormalTok{(}\AttributeTok{size =} \DecValTok{6}\NormalTok{)}
\NormalTok{  ) }\SpecialCharTok{+}
  \FunctionTok{labs}\NormalTok{(}\AttributeTok{x =} \StringTok{"Age Group"}\NormalTok{, }\AttributeTok{y =} \StringTok{"Standardized Relative Abundance"}\NormalTok{, }\AttributeTok{color =} \StringTok{"Phylum"}\NormalTok{)}

\FunctionTok{print}\NormalTok{(p\_phylum)}
\end{Highlighting}
\end{Shaded}

\pandocbounded{\includegraphics[keepaspectratio]{5_age_analyses_final_improved_files/figure-latex/unnamed-chunk-5-2.pdf}}

\begin{Shaded}
\begin{Highlighting}[]
\FunctionTok{cat}\NormalTok{(}\StringTok{"Generated parallel coordinate plot colored by phylum.}\SpecialCharTok{\textbackslash{}n}\StringTok{"}\NormalTok{)}
\end{Highlighting}
\end{Shaded}

\begin{verbatim}
## Generated parallel coordinate plot colored by phylum.
\end{verbatim}

\begin{Shaded}
\begin{Highlighting}[]
\CommentTok{\# Step 14: Save both plots as PDFs}
\CommentTok{\# Save plot colored by cluster}
\FunctionTok{save\_pdf\_figure}\NormalTok{(}\StringTok{"15.1\_parallel\_coord\_plot\_cluster\_improved.pdf"}\NormalTok{, }\AttributeTok{width\_mm =} \DecValTok{55}\NormalTok{, }\AttributeTok{height\_mm =} \DecValTok{45}\NormalTok{)}
\FunctionTok{print}\NormalTok{(p\_cluster)}
\FunctionTok{dev.off}\NormalTok{()}
\end{Highlighting}
\end{Shaded}

\begin{verbatim}
## pdf 
##   2
\end{verbatim}

\begin{Shaded}
\begin{Highlighting}[]
\FunctionTok{cat}\NormalTok{(}\StringTok{"Saved parallel coordinate plot colored by cluster as PDF.}\SpecialCharTok{\textbackslash{}n}\StringTok{"}\NormalTok{)}
\end{Highlighting}
\end{Shaded}

\begin{verbatim}
## Saved parallel coordinate plot colored by cluster as PDF.
\end{verbatim}

\begin{Shaded}
\begin{Highlighting}[]
\CommentTok{\# Save plot colored by phylum}
\FunctionTok{save\_pdf\_figure}\NormalTok{(}\StringTok{"15.2\_parallel\_coord\_plot\_phylum\_improved.pdf"}\NormalTok{, }\AttributeTok{width\_mm =} \DecValTok{75}\NormalTok{, }\AttributeTok{height\_mm =} \DecValTok{45}\NormalTok{)}
\FunctionTok{print}\NormalTok{(p\_phylum)}
\FunctionTok{dev.off}\NormalTok{()}
\end{Highlighting}
\end{Shaded}

\begin{verbatim}
## pdf 
##   2
\end{verbatim}

\begin{Shaded}
\begin{Highlighting}[]
\FunctionTok{cat}\NormalTok{(}\StringTok{"Saved parallel coordinate plot colored by phylum as PDF.}\SpecialCharTok{\textbackslash{}n}\StringTok{"}\NormalTok{)}
\end{Highlighting}
\end{Shaded}

\begin{verbatim}
## Saved parallel coordinate plot colored by phylum as PDF.
\end{verbatim}

\subsubsection{Explanation of the Improved Multi-Method Clustering
Approach}\label{explanation-of-the-improved-multi-method-clustering-approach}

\paragraph{\texorpdfstring{\textbf{Data
Preparation}}{Data Preparation}}\label{data-preparation}

\begin{itemize}
\tightlist
\item
  \textbf{Alpha Diversity}: Loaded from \texttt{dada2rsv-alpha.txt}
  (e.g., Observed OTUs) and merged with sample metadata containing age.
\item
  \textbf{Beta Diversity}: Unweighted UniFrac distances loaded from
  \texttt{dada2rsv-beta-unwunifrac.txt}.
\item
  \textbf{Age Groups}: Individuals are binned into groups (14-19, 20-24,
  \ldots, 80-85) for analysis.
\end{itemize}

\paragraph{\texorpdfstring{\textbf{Simple Clustering
Methodology}}{Simple Clustering Methodology}}\label{simple-clustering-methodology}

The new approach uses a \textbf{simple silhouette method} with minimal
constraints to naturally find 3-6 clusters:

\begin{enumerate}
\def\labelenumi{\arabic{enumi}.}
\tightlist
\item
  \textbf{Simple Silhouette Analysis}: Uses silhouette width as the
  primary metric
\item
  \textbf{No Complex Penalties}: Removes all penalty systems for
  simplicity
\item
  \textbf{Simple Constraints}: Caps at 6 clusters maximum, minimum 3
  clusters
\item
  \textbf{Clean Logic}: Just finds the k with highest silhouette score,
  then constrains to 3-6
\end{enumerate}

\paragraph{\texorpdfstring{\textbf{Simple
Constraints}}{Simple Constraints}}\label{simple-constraints}

\begin{itemize}
\tightlist
\item
  \textbf{Minimum 3 Clusters}: Ensures meaningful biological
  interpretation
\item
  \textbf{Maximum 6 Clusters}: Prevents over-clustering
\item
  \textbf{Data-Driven}: Automatically adjusts to data size constraints
\item
  \textbf{Clean}: No complex penalty calculations
\end{itemize}

\paragraph{\texorpdfstring{\textbf{Benefits of Simple
Approach}}{Benefits of Simple Approach}}\label{benefits-of-simple-approach}

\begin{itemize}
\tightlist
\item
  \textbf{🎯 Clean}: Simple logic, easy to understand and interpret
\item
  \textbf{🧬 Biologically Sensible}: Finds 3-6 clusters (perfect range!)
\item
  \textbf{📊 Data-Driven}: Completely automated, no manual intervention
\item
  \textbf{🔄 Simple}: No complex penalty systems to debug
\item
  \textbf{⚖️ Reliable}: Uses standard silhouette method with simple
  constraints
\end{itemize}

This approach should find a more appropriate number of clusters (3-6)
that better reflects the biological structure of your microbiome data!
🚀

\end{document}
